%  $Id: template.tex,v 1.14 2016/07/14 18:28:18 wawood Exp $

\documentclass[]             % options: RDPonly, coveronly, nocover, backcover
{NASA}                       %   plus standard article class options

% Luna, Gemini Pro, 2024-09-20
%
% Here are some alternative titles, reflecting the content and potential focus of your paper:

% Shorter & Punchier:

% Tensor Train Acceleration of Matrix-Free Finite Element Analysis
% High-Performance Low-Order FEM with Tensor Train Element Evaluation
% Scaling Low-Order FEM: A Matrix-Free Tensor Train Approach
% More Specific:

% Tensor Train Element-by-Element~(EBE) Method for Scalable 3D Elasticity
% Performance Analysis of Tensor Train Elasticity in a Matrix-Free FEM Framework
% Overcoming Memory Bottlenecks in Low-Order FEM with Tensor Trains
% Emphasizing Application:

% Efficient Simulation of Large-Scale 3D Solids using Tensor Train FEM
% Accelerating Multiphysics Tetrahedral Models with a Matrix-Free Tensor Train Method
% When choosing, consider:

% Keywords: What terms will researchers in your field be searching for?
% Journal: Some journals prefer concise titles, others more descriptive ones.
% Uniqueness: Does the title stand out and accurately reflect your novel approach?

%  Fill in the following catagories between the braces {}

%\title{Femera Theory Manual}

%\title{Low-order Matrix-free Finite Element %-by-element
%Node-level Performance with Tensor Train Elasticity Evaluation}

\title{Theory and performance of Femera, a matrix-free finite element implementation
with element-by-element response by tensor train analysis}

% title for CMAME
%\title{A tensor train element-by-element method for high-performance matrix-free finite element analysis}

\AuthorAffiliation{David Wagner\\  % list the authors with affiliations
NASA Langley Research Center\\     % for cover and title pages
%  (second author)\\
%  (second author affiliation)
[10pt]
}
\NasaCenter{Langley Research Center\\Hampton, Virginia 23681-2199}
\Type{TM}                      % TM (one technical reviewer), TP (panel review), CR, CP, SP, TT
\Number{XXXXXXXX}              % Report number
\Month{02}                     % two digit number
\Year{2026}                    % four digit number

% Items if Report Documentation Page desired.
%
% \author{} % list all the author names for the RDP
% \SubjectCategory{}           % two digit number
% \Distribution{}              % E.g., Standard
% \LNumber{}                   % Langley L-number
% \SubjectTerms{}              % 4-5 comma separated words
% \Pages{}                     % Document pages +5 for covers & RDP
% \DatesCovered{}              % 10/2000--9/2002
% \ContractNumber{}            % NAS1-12345
% \GrantNumber{}               % NAG1-1234
% \ProgramElementNumber{}
% \ProjectNumber{}             % NCC1-123
% \TaskNumber{}                % Task 123
% \WorkUnitNumber{}            % 123-45-67-89
% \SupplementaryNotes{}
%
\Acknowledgment{This work was supported by the High-performance Computing Incubator at NASA
 Langley Research Center.
 The author would like to acknowledge the Transformational Tools and Technologies Project of the Transformative Aeronautics Concepts Program for support of this work.
%
%TODO email collaborators about using their names.
%
The author thanks Dr. Gabriele Jost for assistance with memory layout and parallelization 
optimization strategies and Daniel Kokron for the initial vector intrinsics implementation. 
%TODO [Additional names] provided valuable feedback on 
%the manuscript and validation of benchmark results against [specific codes].
Dr. Saikumar Yeratapally provided validation of benchmark results against SciFEN~\cite{ref:scifen}.
%
This manuscript benefited from internal peer review at NASA Langley Research Center. The author thanks 
[Dr. Committee Chair], [Dr. Reviewer~1], [Dr. Reviewer~2], %TODO
[Dr. Supervisor], and [Dr. Division Chief Engineer] %TODO
for their careful review and valuable feedback.
%
Computing resources were provided by the 
NASA Langley Research Center K-Cluster, the
NASA Advanced Supercomputing facility, and 
[branch HPC resources if used for this].%TODO
}            % optional back of title page section

%abstract: suggested 200 words
\abstract{Traditional low-order finite element implementations limit model size by
 the amount, and solution time by the bandwidth, of main memory available
 to store and evaluate a large sparse matrix representing the physics of
 the system.
 Rather than building a large linear system response operator (i.e., sparse matrix) as a sum of
 element operators, the element-by-element method builds the system response
 as the sum of element responses.\footnote{\emph{Source code}:
 \texttt{https://github.com/waggyz/femera-mini-app/tree/v0.1}}
 Each element response can be evaluated efficiently as a sequence of small,
 dense tensor operations: a tensor train.
 Tensor train evaluation can significantly improve iterative solver performance
 while using only a small fraction of the memory required by traditional
 element-by-element and sparse matrix-based implementations.
 Lower memory storage and access requirements improve single-node performance
 for three-dimensional elasticity and multiphysics tetrahedron models having
 system sizes ranging over more than five orders of magnitude.
 Larger models benefit significantly from domain over-decomposition, and
 nested parallelism can be used efficiently to solve batches of smaller
 models concurrently on a single compute node.

Detailed performance analyses of 3D models composed of linear, quadratic,
 or cubic tetrahedrons reveal [TODO...] 

Several research use case studies demonstrate how the performance of linear
 and non-linear models compares to popular finite element applications.
%TODO Is the following neede under current guidance?
\footnote{\emph{Disclosure of AI use}:
The author used generative AI, both to improve readability, and to generate initial drafts of portions of this document. The author retains responsibility for the entirety of the text.}

~

TODO \emph{Summarize the analyses and examples presented to show why this is the best
 thing since sliced bread.}

\emph{FIXME "traditional" is used too many times.}

\emph{FIXME better title}

\emph{TODO consider publishing in: Computer Methods in Applied Mechanics and Engineering (CMAME) [Elsevier]}
\emph{See LaTeX comments for more suggestions.}

~

}
% UPDATE: Computer Methods in Applied Mechanics and Engineering (CMAME)
% Perfect fit - This is likely your best target
% Strong focus on novel numerical methods in FEM
% Emphasizes computational efficiency and implementation
% High impact in the community you're targeting
%
% Please suggest appropriate peer-reviewed journals for this paper.
% LunAI, Google Gemini 1.5 Pro
% Sep 19, 2024, 4:09:38 PM

% WAS: Integrating Materials and Manufacturing Innovation, Tools, Methods, and Impact of Digitally Enabled Materials Engineering, https://link.springer.com/journal/40192}
% Based on the focus on parallel finite element methods, here are some peer-reviewed journals that could be a good fit:

% Top Tier Journals (High Impact & Rigor):

% Journal of Computational Physics: A classic choice for numerical methods in physics, including finite element techniques.
% Computer Methods in Applied Mechanics and Engineering~(CMAME): Excellent for papers with a strong engineering applications focus alongside computational methods.
% International Journal for Numerical Methods in Engineering~(IJNME): Well-respected journal specifically for numerical techniques in various engineering fields.
% Mid-Tier Journals (Still Reputable, Potentially Faster Review):

% Computers & Structures: A good option if your work has applications in structural mechanics or related areas.
% Finite Elements in Analysis and Design: More specialized to finite element methods, suitable if your contributions are primarily methodological.
% Numerical Algorithms: A broader journal, but suitable if your paper has a strong emphasis on novel algorithms within the parallel finite element context.
% Journal of Scientific Computing: A good choice if your work has broad applicability in scientific computing beyond just a single domain.
% Factors to Consider When Choosing:

% Scope and Audience: Make sure the journal's scope aligns with your paper's focus (e.g., engineering applications, specific physics problems, general numerical methods).
% Impact Factor: While not the only metric, it reflects the journal's visibility and potential citation impact.
% Review Time: Some journals have faster review processes than others, which might be important depending on your timeline.
% Open Access Options: If you want your work to be freely available, check the journal's open access policies and fees.
% Additional Tips:

% Read recent publications: Browse through recent issues of potential journals to see if your paper aligns with the type of research they publish.
% Check author guidelines: Each journal has specific formatting and submission guidelines, so make sure to review those carefully.
% I hope this helps you find the perfect journal for your work!

%\usepackage{graphicx}

\usepackage{amsmath}
\usepackage{amssymb}
\usepackage{mathtools}
\usepackage{booktabs}
\usepackage[ruled,vlined]{algorithm2e}
\usepackage[all]{nowidow}
\usepackage{siunitx}

\usepackage{multicol}
\usepackage{ragged2e}
\usepackage{makecell}

\newcommand{\mcone}[2]{\multicolumn{1}{#1}{#2}}

\begin{document}

\newpage\tableofcontents\newpage

\section{Nomenclature}
%\begin{multicols*}{2}

\begin{description}
  \item[$n_\mathrm{p}$] Number of partitions
  \item[$n_\mathrm{\ell}$] Number of elements
  \item[$n_\mathrm{n}$] Number of nodes
  \item[$n_\mathrm{u}$] Number of degrees of freedom~(dof)
  \item[$n_\mathrm{\Lambda}$] Number of integration points
  \item[$d_\mathrm{m}$] Model spatial dimension
  \item[$d_\mathrm{\ell}$] Element spatial dimension
  \item[$d_\mathrm{u}$] Node dimension (degrees of freedom per node)
  \item[$\mathbb{P}$] Set of all partitions of the model
  \item[$\mathcal{P}$] Set of partitions assigned to a CPU, one ($\mathcal{P}in\mathbb{P}$) for each CPU 
  \item[$\ell$] Element index
  \item[$\Lambda$] Integration point index
  \item[$\mathbf{f}_\ell, {f}_\ell$] Element force vector, second-order tensor
  \item[$\mathbf{K}_\ell, {K}_\ell$] Element stiffness matrix, fourth-order tensor
  \item[$\mathbf{u}_\ell, {u}_\ell$] Element displacement vector, second-order tensor
  \item[$\mathbf{B}_\Lambda$] Strain-displacement matrix at integration point~$\Lambda$
  \item[$\mathbf{G}_\Lambda, G_\Lambda$] Geometric distortion matrix, second-order tensor, at integration point~$\Lambda$
  \item[$\mathbf{J}_\Lambda, J_\Lambda$] Geometric Jacobian matrix, second order tensor, at integration point~$\Lambda$
  \item[$\dot{\mathbf{N}}_{\Lambda}, \dot{N}_{\Lambda}$] Shape function derivative matrix, second-order tensor, at integration point~$\Lambda$
  \item[$\mathbf{D}, C$] Material property (Voigt) matrix in global coordinates, fourth-order tensor
  \item[$\widetilde{\mathbf{D}}, \widetilde{C}$] Material property (Voigt) matrix in material coordinates, fourth-order tensor

%  \item[$\varepsilon_\text{V}$]
  \item[$\vec{\varepsilon}, \epsilon$] Engineering strain Voigt vector\cite{ref:voigt1928lehrbuch,ref:love2013treatise} in global coordinates, true strain second-order tensor
  \item[$\widetilde{\epsilon}$] True strain second-order tensor in material coordinates
%  \item[$\sigma_\text{V}$]
  \item[$\vec{\sigma}, \sigma$] Cauchy stress Voigt vector in global coordinates, second-order tensor
  \item[$\widetilde{\sigma}$] Cauchy stress second-order tensor in material coordinates
  
  \item[$\mathbf{Q}, \mathbf{Q}'$] Material orientation transformation matrices
%  \item[$\epsilon$] True strain second-order tensor
%  \item[$\sigma$] Cauchy stress second-order tensor
  \item[$R$] Material orientation second-order tensor
  \item[$M_{\mathcal{G}}$] Preconditioner
  \item[$r$] Residuals for each degree of freedom
  \item[$p$] Conjugate gradient descent direction
  \item[$\zeta^{2}$] Preconditioned residual norm squared
  \item[$\zeta_{\mathrm{tol}}^{2}$] Preconditioned residual tolerance squared
  \item[$\alpha$] Conjugate gradient step size

\end{description}
%\end{multicols*}

\newpage\listoftables\newpage
\listoffigures\newpage\listofalgorithms\newpage

% Hero's journey (Vogler)
% https://en.wikipedia.org/wiki/The_Writer%27s_Journey:_Mythic_Structure_for_Writers
%
% The archetypes
%
% Hero:               someone who is willing to sacrifice his own needs on behalf of others
%                     * FEM (Algebra from linear algebra to tensor algebra)
% Mentor:             all the characters who teach and protect heroes and give them gifts
%                     * Continuum mechanics
%                     * Einstein convention
% Threshold Guardian: a menacing face to the hero, but if understood, they can be overcome
%                     * Traditional FE presentation, legacy code
% Herald:             a force that brings a new challenge to the hero
%                     * Performance optimization for modern processors
% Shapeshifter:       characters who change constantly from the hero's point of view
%                     * Constitutive matrix/tensor/operator
% Shadow:             character who represents the energy of the dark side
%                     * Arithmetic intensity, memory latency and bandwidth, interprocess communication
%                     * Mechanics of materials, linear algebra
% Ally:               someone who travels with the hero through the journey, serving variety of functions
%                     * Performance optimization techniques
%                     * Continuum mechanics, Einstein notation
% Trickster:          embodies the energies of mischief and desire for change
%                     * high performance architecture, processor technology
%                     * very large problems (microstructure, complex geometry), optimization, UQ, inverse problems
%
% Stages of the Journey
%
%   I. Departure =============================================================
%HJ 1. Ordinary world --------------------------------------------------------
%      the hero is seen in their everyday life
%HJ 2. Call to adventure -----------------------------------------------------
%      the initiating incident of the story
%HJ 3. Refusal of the call ---------------------------------------------------
%      the hero experiences some hesitation to answer the call
%HJ 4. Meeting with the mentor -----------------------------------------------
%      the hero gains the supplies, knowledge, and confidence needed to 
%      commence the adventure
%HJ 5. Crossing the first threshold ------------------------------------------
%      the hero commits wholeheartedly to the adventure
%  II. Initiation ============================================================
%HJ 6. Tests, allies, and enemies --------------------------------------------
%      the hero explores the special world, faces trials, and makes friends and 
%      enemies
%HJ 7. Approach to the inmost cave -------------------------------------------
%      the hero nears the center of the story and the special world
%HJ 8. The ordeal ------------------------------------------------------------
%      the hero faces the greatest challenge yet and experiences death and 
%      rebirth
%HJ 9. Reward ----------------------------------------------------------------
%      the hero experiences the consequences of surviving death
% III, Return ================================================================
%HJ10. The road back ---------------------------------------------------------
%      the hero returns to the ordinary world or continues to an ultimate 
%      destination
%HJ11. The resurrection ------------------------------------------------------
%      the hero experiences a final moment of death and rebirth so they are 
%      pure when they reenter the ordinary world
%HJ12. Return with the elixir ------------------------------------------------
%      the hero returns with something to improve the ordinary world

\section{Introduction}% Original text by Luna, Gemini 1.5 Pro, 2024-09-08 
\label{sec:Introduction}
%HJ 1. Ordinary world *******************************************************
The finite element method~(FEM) is a powerful numerical technique for obtaining approximate solutions to partial differential equations~(PDEs), making it an indispensable tool in many engineering disciplines for analyzing and designing complex physical systems.  
Low-order (first-, second-, and third-order) finite element implementations are particularly well-suited for analyzing implicit, static, or quasi-static equilibrium behavior of solids and structures.
%HJ 2. Call to adventure ****************************************************
% TODO need this step.

%HJ 3. Refusal of the call **************************************************
However, traditional low-order finite element analysis frequently encounters computational bottlenecks, particularly when addressing large-scale problems. 
These bottlenecks stem from the storage and evaluation of a large, sparse matrix representing the linear response of the system. 
The size of these matrices place limitations on the scale of models that can be effectively handled. 
Furthermore, computation of the system response is inefficient, limited by main memory bandwidth, because of the low arithmetic intensity of sparse matrix-vector~(SpMV) multiplication. 

%HJ 4. Meeting with the mentor **********************************************
%TODO Move this to first paragraph?
A recent finite element matrix-free approach~\cite{ref:settgast2023performant} 
overcomes these limitations and demonstrates efficient computational performance when the code and data needed for the linear solve fit entirely within the memory of a single graphics processing unit~(GPU) card.
This work demonstrates that this approach, herein referred to as element-by element~(EBE) response tensor train analysis~(EBERTTA), is also efficient for single-node central processing unit~(CPU) evaluation. %, as demonstrated in section~\ref{sec:Examples}.
%HJ 5. Crossing the first threshold *****************************************
The mathematical basis for EBERTTA can be derived by recasting finite element operations from the traditional linear algebra presentation % in section~\ref{subsec:Matrix-Assembly}
to tensor algebra. %, as described in sections~\ref{sec:Methods} and~\ref{sec:Methods}.
Implementation of these tensor operations presents several opportunities for CPU performance optimization.

EBERTTA represents each element response as a tensor train, a term borrowed from the field of machine learning.\cite{ref:oseledets-2011-tensor-train} 
The EBERTTA approach circumvents the need to assemble the global stiffness matrix, instead constructing the global response by summing the contributions from each element's response. 
Representing each element response as a tensor train allows efficient evaluation through a sequence of small, dense tensor operations. 
This approach significantly reduces the memory footprint, thereby accommodating larger models, and enables more efficient memory access, leading to enhanced solver performance.

%HJ Foreshadowing ***********************************************************
Detailed performance analyses on three-dimensional elasticity models having linear, quadratic, or cubic tetrahedral elements quantify the efficacy of EBERTTA. 
These analyses encompass a wide range of system sizes spanning five orders of magnitude and highlight the performance benefits across different CPU hardware architectures. 
Furthermore, several case studies illustrate the applicability and effectiveness of this approach for both linear and non-linear material models. These studies showcase the competitive performance of the EBERTTA method against established finite element analysis software packages.  

Subsequent sections of this paper delve into the development, implementation, and performance testing of this method. 
%  II. Initiation ============================================================
%HJ 6. Tests, allies, and enemies *******************************************
Section~\ref{sec:Background} provides an overview of traditional finite element iterative solution methods and the challenges associated with large-scale problems. 
%HJ 7. Approach to the inmost cave ******************************************
Section~\ref{sec:Development} details the development of the tensor train response evaluation method, which forms the cornerstone of the EBERTTA approach. 
Section~\ref{sec:Methods} presents the tensor train response evaluation in detail. 
%HJ 8. The ordeal ***********************************************************
Section~\ref{sec:Examples} provides a comprehensive performance analysis of the tensor-train-based EBE method,  contrasting it with existing methods and evaluating its single-node performance characteristics using various  benchmark problems. %TODO the examples are not  really benchmark problems
%HJ 9. Reward ***************************************************************
%TODO brag a little here 
%
%Finally, section~\ref{sec:Discussion} summarizes the findings and their implications, 
%suggesting avenues for future research in section~\ref{sec:Future-Work}.
% III, Return ================================================================
%HJ10. The road back ********************************************************
Finally, section~\ref{sec:Discussion} sums up the results that 
lead to the conclusions of section~\ref{sec:Conclusions}, and 
%HJ11. The resurrection *****************************************************
point the way toward further research outlined in section~\ref{sec:Future-Work}.
%HJ12. Return with the elixir ***********************************************

\section{Background}
\label{sec:Background}
%*****************************************************************************
%HJ 1. Ordinary world
%
%
% Intro paragraph draft by Luna Gemini Pro, 2024-09-15
Solving large sparse linear systems arising from discretized partial differential equations often necessitates iterative methods. Direct solvers, while providing solutions up to the limit of floating point accuracy, become computationally intractable for the problem sizes considered in this work. Iterative solvers, such as the preconditioned conjugate gradient~(PCG) method\cite{ref:barrett-1994}, offer an efficient alternative by successively refining an initial guess until convergence within a specified tolerance.

% Orignal draft by Luna Gemini Pro, 2024-09-24
While iterative methods offer a path to solving large-scale linear systems, their efficiency often hinges on parallelization techniques. Domain decomposition is one such technique that addresses this challenge by dividing the computational load.

\subsection{Domain decomposition}

% Draft by Luna Gemini Pro, 2024-09-24, edited by the author
Domain decomposition is a technique for parallelizing the solution of a finite element model by dividing it into smaller, more manageable subdomains. Each subdomain can then be assigned to a different processor or core, enabling parallel computation. This approach is essential for tackling large-scale problems that would be computationally prohibitive for a single processor to handle.

This work considers domain decomposition methods that divide a model along element faces into non-overlapping partitions. 
%TODO define halo nodes and elements
Halo  exchange, a crucial aspect of domain decomposition, involves communication between processors to ensure consistency of data along the boundaries of adjacent partitions. This exchange typically involves synchronizing values associated with shared nodes (referred to as halo, rind, or ghost nodes) between neighboring partitions.

Domain decomposition introduces performance implications due to the need for non-local memory access and atomic operations. When a processor needs to access data associated with halo nodes residing in another processor's memory, it incurs the overhead of inter-processor communication, which can be orders of magnitude slower than local memory access. Additionally, updating shared data often requires atomic operations to prevent race conditions, where multiple processors attempt to modify the same memory location simultaneously. Atomic operations introduce synchronization overhead, potentially hindering parallel efficiency.

The choice of domain decomposition strategy and the implementation of halo exchange significantly influence the overall performance of parallel finite element analysis. 
% original by auther
In this work, models are decomposed into partitions, each having approximately the same number of elements. All of each partition's halo nodes are present in the partition, essentially duplicating those nodes shared along the boundaries of partitions, with a corresponding increase the memory needed for the model. One of the duplicate shared nodes is assigned as the primary node to gather the contributions of the duplicate nodes before scattering the accumulated results back to the contributing partitions.

Despite increased memory needs, this strategy can improve performance because computational kernels access memory that is closer (in terms of memory latency) to each processor. This domain decomposition approach can also take advantage of memory latency hiding\cite{ref:latency-hiding} to improve performance.

\subsection{Hiding parallel synchronization latency}
\label{subsec:hide-parallel-sync-latency}

% Intro paragraph by Luna Gemini Pro, 2024-09-15, edited by the author
Parallel implementations of iterative solvers often face performance bottlenecks due to communication overhead during halo exchanges. This issue can be mitigated by overlapping communication with computation. Specifically, processors can perform computations on their local partitions while waiting for halo data synchronization, effectively hiding some of the latency associated with synchronization.

The domain decomposition strategy in section~\ref{subsec:Domain-Decomposition} allows for a three-phase approach to compute the values associated with the nodes in each partition.
\begin{enumerate}
\item{Phase one evaluates halo nodes or elements containing halo nodes.}
\item{Phase two does the halo exchange.}
\item{Phase three evaluates interior nodes or elements.}
\end{enumerate}
Non-halo interior nodes and elements do not have any halo nodes and therefor do not participate in synchronization. This approach allows for some of the time each CPU needs for atomic operations and non-local memory access in phase two to overlap the computational evaluation of interior nodes that access only non-shared local memory in phase three.

% Transition draft by Luna Gemini Pro, 2024-09-27
Beyond parallelization, the convergence rate of iterative solvers can be significantly enhanced by employing preconditioners.

\subsection{Jacobi preconditioner}
\label{subsec:jacobi-preconditioner}

% Draft by Luna Gemini Pro, 2024-09-15
This work utilizes a Jacobi preconditioner, a simple yet effective method for improving the convergence rate of iterative solvers. The Jacobi preconditioner is a diagonal matrix formed by the inverse of the diagonal elements of the original matrix. Its application involves simple element-wise multiplication, making it computationally inexpensive and straightforward to parallelize.

% Transition by Luna Gemini Pro, 2024-09-27
Combining domain decomposition with a suitable preconditioner forms the foundation for an efficient parallel iterative solver.
%Algorithm~\ref{alg:halo-pcg} outlines the steps involved in a latency-hiding preconditioned conjugate gradient (PCG) method, which leverages these techniques.

\subsection{Iterative solvers}
\label{subsec:Solvers}

% Intro draft by Luna Gemini Pro, 2024-09-15, edited by the author
The steps involved in the latency-hiding preconditioned PCG algorithm~\ref{alg:halo-pcg} are outlined below.
% Algorithm explanation draft by Luna Gemini Pro, 2024-09-15, edited by the author
%Explanation of the Algorithm:

\begin{enumerate}
    \item{\textbf{Initialization}:
    \begin{itemize}
        \item Each processor computes the preconditioner~($M_{\mathcal{G}}$) and initial re\-sponse~(${f_{\mathcal{G}}}$) for its local partition in the set of partitions being processed in parallel.
        \item Synchronize halo nodes, representing data shared among partitions.
        \item Calculate initial residuals~($r$), conjugate gradient descent direction~($p$), and squared preconditioned residual norm~($\zeta^2$).
     \end{itemize}}
    \item{\textbf{Iteration}:
    \begin{itemize}
        \item The algorithm iterates until the preconditioned squared conjugate residual~($\zeta^2$) falls below a squared specified tolerance~($\zeta_{\mathrm{tol}}^{2}$).
        \item \textbf{Halo Exchange}: Each processor computes the conjugate response contributions from halo elements and synchronizes the results.
        \item{\textbf{Kernel PCG Computations}: The algorithm evaluates three main computational kernels (PCG~A, B, and C) within each iteration:
        \begin{itemize}
            \item \textbf{Kernel PCG~A}: Calculates the step size~($\alpha$) by accumulating the dot product of the conjugate direction and residual.
            \item \textbf{Kernel PCG~B}: Updates the residuals, applies the preconditioner, and computes the new preconditioned residual.
            \item \textbf{Kernel PCG~C}: Updates the solution and the conjugate direction.
        \end{itemize}}
    \end{itemize}}
    \item \textbf{Convergence}: The algorithm checks the convergence criterion at the beginning of each iteration.
\end{enumerate}

% Algorithm by the author
\begin{algorithm}[ht] %WAS hbt!
\caption{Latency-hiding parallel preconditioned conjugate gradient method}
\label{alg:halo-pcg}
\begin{enumerate}
\item Initialize:
\begin{enumerate}
\item For each partition in partition set~($\mathcal{P} \subset \mathbb{P}$),
\begin{enumerate}
\item Compute preconditioner: $M_{\mathcal{G}}$.
\item Compute initial response: $f_{\mathcal{G}}=K({u}_{0})$.
\end{enumerate}
\item Synchronize halo node values: $f_{\mathcal{G}},M_{\mathcal{G}}$ (atomic).
\item For each partition in partition set~($\mathcal{P}$),
\begin{enumerate}
\item Invert $M$.
\item Compute initial (negative) residuals: ${r}={b}-{f}$.
\item Compute initial (gradient descent) direction: ${p}=M^{-1}({r})$.
\item Reduce preconditioned residual norm: $\zeta_{P}^2={r}_{\mathcal{P}}\cdot{p}_{\mathcal{P}}$.
\item Reduce $\zeta^{2}=(\zeta^{2}+\sum_{P}\zeta_{P}^2$ (atomic).
\end{enumerate}
\end{enumerate}
\item Iterate: while $\zeta^{2}>\zeta_{\mathrm{tol}}^{2}$ do:
\begin{enumerate}
\item For each partition in partition set~($\mathcal{P}$) 
\newline
Compute halo element conjugate response: $f_{\mathcal{J}}=\sum_{\mathcal{H}}K_{\mathcal{H}}(p)$.
\item Synchronize halo node values: $f_{\mathcal{G}}$ (atomic).
\item For each partition in partition set~($\mathcal{P}$),
\begin{enumerate}
\item Accumulate interior element conjugate response: $f_{\mathcal{J}}=f_{\mathcal{J}}+\sum_{\mathcal{I}}K_{\mathcal{I}}(p)$.
\item Evaluate \textbf{Kernel PCG A}: $\gamma_{P}={p}_{\mathcal{\mathcal{P}}}\cdot{f}_{\mathcal{P}}$.
\item Reduce $\gamma=\gamma+\sum_{P}\gamma_{P}$ (atomic).
\end{enumerate}
\item Calculate $\alpha=\frac{\zeta^{2}}{\gamma}$.
\item For each partition in partition set~($\mathcal{P}$),
\begin{enumerate}
\item Evaluate \textbf{Kernel PCG B}:
\begin{enumerate}
\item Update residuals: ${r}={r}-\alpha{f}$.
\item Apply preconditioner to residuals ${z}=M^{-1}({r})$.
\item Reduce preconditioned residual norm: $\zeta'_{p}={r}_{\mathcal{P}}\cdot{z}_{\mathcal{P}}$.%TODO check this
\end{enumerate}
\item Reduce $\zeta'=\zeta'+\sum_{P}\zeta'_{P}$ (atomic).%TODO check this.
\end{enumerate}
\item Calculate $\beta=\frac{\zeta'}{\zeta^{2}}$.
\item For each partition in partition set~($\mathcal{P}$), evaluate \textbf{Kernel PCG C}:
\begin{enumerate}
\item Update solution: ${u}={u}+\alpha{p}$.
\item Update conjugate direction: ${p}={z}+\beta{p}$.
\end{enumerate}
\item For the next iteration, update: $\zeta^{2}=\zeta'$.
\end{enumerate}
\end{enumerate}
\end{algorithm}

% Original text by author
% A solution~($u$) can be considered converged when the preconditioned squared conjugate
%  residual norm~($\zeta^{2}$) is within a specified tolerance~($\zeta_{\mathrm{tol}}^{2}$),
% \begin{eqnarray}
% \zeta^{2} & \leq & \zeta_{\mathrm{tol}}^{2}\text{.}\label{eq:fem-global-system}
% \end{eqnarray}

% Transition paragraph draft by Luna Gemini Pro, 2024-09-15, edited by the author
Having established the iterative solution strategy, the next two subsections delve into traditional methods of computing the system response~$K(p)$.

\subsection{Sparse matrix assembly and evaluation}
\label{subsec:Matrix-Assembly}

Typical implicit finite element implementations build a system of linear equations,
 representing global nodal force equilibrium conditions, with
 a large, sparse, symmetric, positive-definite stiffness matrix~($\mathbf{K}$)
 acting as a linear response operator.
 
Sparse stiffness matrix assembly starts by computing an element's
Jacobian~($\mathbf{J}_{\Lambda}$) at an integration point~($\Lambda$)
within each element~($\ell$).
 This Jacobian represents how the element has been distorted from
 its natural shape.
 Calculate a 3D~Jacobian as a matrix-matrix product of the shape function derivatives at the integration
 point~($\dot{\mathbf{N}}_{\Lambda}$) and the element's nodal coordinates~($\mathbf{x}_{\ell}$), 
\begin{eqnarray}
\mathbf{J}_{\Lambda} & = & \dot{\mathbf{N}}_{\Lambda}\mathbf{x}_{\ell}\label{eq:geometric-jacobian}\\
\left[\begin{array}{ccc}
\frac{\partial x}{\partial r} & \frac{\partial y}{\partial r} & \frac{\partial z}{\partial r}\\
\frac{\partial x}{\partial s} & \frac{\partial y}{\partial s} & \frac{\partial z}{\partial s}\\
\frac{\partial x}{\partial t} & \frac{\partial y}{\partial t} & \frac{\partial z}{\partial t}
\end{array}\right] & = & \left[\begin{array}{ccc}
\frac{\partial N_{1}}{\partial r} &  & \frac{\partial N_{n}}{\partial r}\\
\frac{\partial N_{1}}{\partial s} & \cdots & \frac{\partial N_{n}}{\partial s}\\
\frac{\partial N_{1}}{\partial t} &  & \frac{\partial N_{n}}{\partial t}
\end{array}\right]\left[\begin{array}{ccc}
x_{1} & y_{1} & z_{1}\\
 & \vdots\\
x_{n} & y_{n} & z_{n}
\end{array}\right]\text{.}\nonumber 
\end{eqnarray}
 Apply the inverse of the Jacobian to the shape function natural coordinate
 gradient to calculate the shape function gradient in global coordinates~($\mathbf{G}_{\Lambda}$),
\begin{eqnarray}
\mathbf{G}_{\Lambda} & = & \mathbf{J}_{\Lambda}^{-1}\dot{\mathbf{N}}_{\Lambda}\label{eq:shape-gradient-global}\\
\left[\begin{array}{ccc}
\frac{\partial N_{1}}{\partial x} &  & \frac{\partial N_{n}}{\partial x}\\
\frac{\partial N_{1}}{\partial y} & \cdots & \frac{\partial N_{n}}{\partial y}\\
\frac{\partial N_{1}}{\partial z} &  & \frac{\partial N_{n}}{\partial z}
\end{array}\right] & = & \left[\begin{array}{ccc}
\frac{\partial r}{\partial x} & \frac{\partial s}{\partial x} & \frac{\partial t}{\partial x}\\
\frac{\partial r}{\partial y} & \frac{\partial s}{\partial y} & \frac{\partial t}{\partial y}\\
\frac{\partial r}{\partial z} & \frac{\partial s}{\partial z} & \frac{\partial t}{\partial z}
\end{array}\right]^{-1}\left[\begin{array}{ccc}
\frac{\partial N_{1}}{\partial r} &  & \frac{\partial N_{n}}{\partial r}\\
\frac{\partial N_{1}}{\partial s} & \cdots & \frac{\partial N_{n}}{\partial s}\\
\frac{\partial N_{1}}{\partial t} &  & \frac{\partial N_{n}}{\partial t}
\end{array}\right]\text{.} \nonumber
\end{eqnarray}

 A sparse strain-displacement matrix~($\mathbf{B}_{\Lambda}$) 
 containing the elements of~($\mathbf{G}_{\Lambda}$) is a linear operator that computes a Voigt
 vector representing the engineering strain~($\vec{\varepsilon}_{\Lambda}$)
 at an integration point ($\Lambda$) when it operates on an unrolled
 element displacement vector~($\mathbf{u}_{\ell}$),
\begin{eqnarray}
\vec{\varepsilon}_{\Lambda} & = & \mathbf{B}_{\Lambda}\mathbf{G}_{\Lambda}\mathbf{u}_{\ell}\label{eq:voigt-strain-from-displacement}\\
\left[\begin{array}{c}
\varepsilon_{x}\\
\varepsilon_{y}\\
\varepsilon_{z}\\
\gamma_{yz}\\
\gamma_{xz}\\
\gamma_{xy}
\end{array}\right]=\left[\begin{array}{c}
\frac{\partial u_{x}}{\partial x}\\
\frac{\partial u_{y}}{\partial y}\\
\frac{\partial u_{z}}{\partial z}\\
\frac{\partial u_{y}}{\partial z}+\frac{\partial u_{z}}{\partial y}\\
\frac{\partial u_{x}}{\partial z}+\frac{\partial u_{z}}{\partial x}\\
\frac{\partial u_{x}}{\partial y}+\frac{\partial u_{y}}{\partial x}
\end{array}\right] 
& = & \left[\begin{array}{ccccccc}
\frac{\partial N_{1}}{\partial x} & 0 & 0 &  & \frac{\partial N_{n}}{\partial x} & 0 & 0\\
0 & \frac{\partial N_{1}}{\partial y} & 0 &  & 0 & \frac{\partial N_{n}}{\partial y} & 0\\
0 & 0 & \frac{\partial N_{1}}{\partial z} & \ldots & 0 & 0 & \frac{\partial N_{n}}{\partial z}\\
0 & \frac{\partial N_{1}}{\partial z} & \frac{\partial N_{1}}{\partial y} &  & 0 & \frac{\partial N_{n}}{\partial z} & \frac{\partial N_{n}}{\partial y}\\
\frac{\partial N_{1}}{\partial z} & 0 & \frac{\partial N_{1}}{\partial x} &  & \frac{\partial N_{n}}{\partial z} & 0 & \frac{\partial N_{n}}{\partial x}\\
\frac{\partial N_{1}}{\partial y} & \frac{\partial N_{1}}{\partial x} & 0 &  & \frac{\partial N_{n}}{\partial y} & \frac{\partial N_{n}}{\partial x} & 0\\
\end{array}\right]\left[\begin{array}{c}
u_{1x}\\
u_{1y}\\
u_{1z}\\
\vdots\\
u_{nx}\\
u_{ny}\\
u_{nz}
\end{array}\right]\text{.} \nonumber 
\end{eqnarray}

Three-dimensional elastic material response is represented by a symmetric~($6\times6$)
Voigt matrix~($\widetilde{\mathbf{D}}$) in material coordinates.
 Homogeneous isotropic material response is invariant with respect to material orientation,
 and the response in global coordinates is the same as it is in material
 coordinates~($\mathbf{D}_{\ell}=\widetilde{\mathbf{D}}$).
 For homogeneous anisotropic materials, the material orientation with respect to the
 global coordinate system can be applied using ($6\times6$)~matrices
 ($\mathbf{Q}_{\ell}$) and ($\mathbf{Q}_{\ell}^{'}$).
 These matrices are constructed such that their entries correspond to the
 Voigt vector representations of engineering strains~($\vec{\varepsilon}$)
 and Cauchy stresses~($\vec{\sigma}$), respectively.
 Material orientation is typically applied as,
\begin{eqnarray}
\mathbf{D}_{\ell} = \mathbf{Q}_{\ell}^{'}\widetilde{\mathbf{D}}\mathbf{Q}_{\ell}\text{,}\label{eq:material-voigt-rotate}
\end{eqnarray}
to transform the material response to the global coordinate system.
 The material response in the global system can then be used, for example,
 to compute a Voigt vector of Cauchy stresses~($\vec{\sigma}$) from a Voigt vector of
 engineering strains~($\vec{\varepsilon}$) as,
\begin{eqnarray}
\vec{\sigma}_{\Lambda} = \mathbf{D}\vec{\varepsilon}_{\Lambda}\text{.}\label{eq:stress-from-engineering-strain-global}
\end{eqnarray}
A dense symmetric element stiffness matrix can be evaluated as,
\begin{eqnarray}
\label{eq:element-stiffness-matrix-traditional}
\mathbf{K}_{\ell} = \sum_{\Lambda}\omega_{\Lambda}|\mathbf{J}_{\Lambda}|\mathbf{B}_{\Lambda}^\mathrm{T}\mathbf{D}_{\ell}\mathbf{B}_{\Lambda}\text{.}\label{eq:la-elem-stiff}
\end{eqnarray}
Summing the entries of each local element stiffness matrix~($\mathbf{K}_{\ell}$)
into their appropriate locations within~($\mathbf{K}$), assembles the global system matrix as, 
\begin{eqnarray}
\mathbf{K} = \sum_{\ell}\mathbf{K}_{\ell}\text{.}\label{eq:global-stiff-assemble}
\end{eqnarray}
A sparse matrix-vector~(SpMV) multiplication computes the system response from an input vector~($\mathbf{u}$) as,
\begin{eqnarray}
\mathbf{f} = \mathbf{K}\mathbf{u}\text{.}
\end{eqnarray}

After assembly of the sparse system matrix, a physics response evaluation
 reads the input non-uniformly, while the compressed sparse system matrix
 and output data can be accessed uniformly, as described in Algorithm~\ref{alg:elastic-spmv}.
 The sparse matrix row~(SpMR) computation kernel is a dense vector dot product that reduces the nonzero entries in each row of the stiffness matrix and the corresponding entries in an input vector.
\begin{algorithm}[ht]%WAS hbt!
\caption{Linear response SpMV~physics evaluation}
\label{alg:elastic-spmv}
\begin{enumerate}
\item Initialize:
\begin{enumerate}
\item Index sets ($\mathcal{C}_\mathcal{R}$) for the nonzero elements in each row ($\mathcal{R}$).
%\item Sparse stiffness matrix: $\mathbf{K}=\sum_{\ell}\mathbf{K}_{\ell}$.
\item Compressed sparse row stiffness matrix: $\mathbf{K_{\mathrm{CSR}}} \Leftarrow \sum_{\ell}\mathbf{K}_{\ell}$.
%\mathbf{K}_\mathcal{R}(\mathbf{K},\mathcal{C}_\mathcal{R}
\end{enumerate}
\item For each partition~($\mathrm{p} \subset \mathbb{P}$):
\begin{enumerate}
\item For each row~(${\mathcal{R}}\in\mathbf{K_{\mathrm{CSR}}}$):
\begin{enumerate}
\item Read: $\mathbf{K}_\mathcal{R}(\mathbf{K_{\mathrm{CSR}}},\mathcal{C}_\mathcal{R})$,
$\mathbf{f_\mathcal{R}(\mathbf{f},\mathcal{C}_\mathcal{R})}$,
$\mathbf{u}_\mathcal{R}(\mathbf{u},\mathcal{C}_\mathcal{R})$.
\item Reduce this response component with a dense dot product:
\newline
\textbf{Kernel SpMR}: $\mathbf{f}_{\mathcal{R}}=\mathbf{K}_{\mathcal{R}}\cdot{\mathbf{u}}_{\mathcal{R}}$.
\item Write: $\mathbf{f}_{\mathcal{R}}$ into ${\mathbf{f}}$.
\end{enumerate}
\end{enumerate}
\item Synchronize $\mathbf{f}$.
\end{enumerate}
\end{algorithm}

%HJ 2. Call to adventure ****************************************************
Iterative Krylov subspace solution methods, such as the conjugate gradient
 method summarized in section~\ref{subsec:Solvers}, do not require the system matrix to be applied directly.
 These methods only need the response of the system to one or more inputs.

% Transition paragraph draft by Luna Gemini Pro, 2024-09-27
%TODO reword this
A critical aspect of iterative solvers in the context of FEM is the computation of the system response, often involving the stiffness matrix. Traditional approaches typically employ either sparse matrix assembly or element-by-element methods. The latter, detailed below, bypasses the formation of a sparse global stiffness matrix.

\subsection{Element-by-element local matrix storage}
\label{subsec:EBE-EMS}

Implementations of the traditional element-by-element method do not assemble the global
 stiffness matrix.
 Instead, each element's stiffness is computed by equation~\ref{eq:la-elem-stiff}
and stored independently, without assembling the global stiffness matrix.
 The global linear response is assembled by summing the contributions
 from each element's response as,
\begin{eqnarray}
\label{ebe-ems-force-summation}
\mathbf{f} & = & \sum_{\ell}\mathbf{f}_{\ell}(\mathbf{u})\text{, with}\label{eq:EBE-force-assembly} \\
\mathbf{f}_{\ell}(u) & = & \mathbf{K}_{\ell}\mathbf{u}_{\ell}\text{.}\label{eq:EMS-elem-response}\text{.} \nonumber
\end{eqnarray}
Traditional element matrix storage~(EMS) EBE algorithms compute each element's
 response by a dense matrix-vector product using symmetric element stiffness
 matrices ($K_{\ell}$) that are stored in memory.
%($K_{\ell}$),
%\begin{eqnarray}
%f_{\ell}(u) & = & \mathbf{K}_{\ell}{u}_{\ell}\text{.}\label{eq:EMS-elem-response}
%\end{eqnarray}

Algorithm~\ref{alg:elastic-tet-lma} summarizes the steps in an EBE EMS linear response evaluation.
 The computational kernel~(EMS\nobreakdash-F) is a dense symmetric matrix-vector product.
\begin{algorithm}[ht]%WAS hbt!
\caption{Linear response EBE EMS physics evaluation}
\label{alg:elastic-tet-lma}
\begin{enumerate}
\item Initialize:
%\begin{enumerate}
%\item E
element stiffness matrices ($\mathbf{K}_{\ell}$)
%\end{enumerate}
\item For each partition set~($\mathcal{P_{s}} \subset P$),
\begin{enumerate}
\item For each element,
\begin{enumerate}
\item Read: $\mathbf{u}_{\ell}$,$\mathbf{f}_{\ell}$, $\mathbf{K}_{\ell}$.
\item Accumulate nodal forces as
\newline
\textbf{Kernel EBE\nobreakdash-EMS F}: $\mathbf{f}_{\ell}=\mathbf{f}_{\ell}+\mathbf{K}_{\ell}\mathbf{u}_{\ell}$.
\item Write: $\mathbf{f}_{\ell}$ into $\mathbf{f}$.
\end{enumerate}
\end{enumerate}
\item Synchronize $\mathbf{f}$.
\end{enumerate}
\end{algorithm}
This approach accesses both the input~($\mathbf{u}_{\ell}$)
and output~($\mathbf{f}_{\ell})$ data non-uniformly,
with the element stiffness matrices~($\mathbf{K}_{\ell}$)
stored so that they can be read uniformly during evaluation.
 
The traditional approaches summarized in the last two sections amortize the computational
cost of one-time evaluation of the element operations of sections~\ref{subsec:Matrix-Assembly}
or \ref{subsec:EBE-EMS}
over the much larger linear system solve time.
 More recent element-by-element approaches recompute element response matrices, during
 each iteration.

\subsection{Element-by-element local matrix construction}
\label{subsec:EBE-LMC}

A more recent method, element-by-element local matrix construction (LMC), does not store element matrices during initialization. Instead, each element matrix is reconstructed by equation~\ref{eq:element-stiffness-matrix-traditional}, then evaluated to compute the element's 
contribution to the global response during each iteration with equation~\ref{eq:EBE-force-assembly}. 
Algorithm~\ref{alg:elastic-tet-lmc} summarizes the steps in an EBE LMC linear response evaluation.
 The computational kernel~(LMC\nobreakdash-K) represents a series of matrix-matrix products involving sparse % $\mathbf{B}$
 B\nobreakdash-matrices, typically implemented as a loop nest.
 The computational kernel~(LMC\nobreakdash-F) is a dense symmetric matrix-vector product.
\begin{algorithm}[ht]%WAS hbt!
\caption{Linear response EBE LMC physics evaluation}
\label{alg:elastic-tet-lmc}
\begin{enumerate}
\item Initialize: Store Jacobian determinanst ($|\mathbf{J}_{ell}|$) and inverse Jacobians ($\mathbf{J}_{\ell}^{-1}$) for each element. 
\item For each partition set~($\mathcal{P_{s}} \subset P$):
\begin{enumerate}
\item For each element,
\begin{enumerate}
\item Read: $\mathbf{u}_{\ell}$,$\mathbf{f}_{\ell}$, $|\mathbf{J}_{ell}|$, $\mathbf{J}_{\ell}^{-1}$.
\item Construct:element stiffness matrix
\newline
\textbf{Kernel LMC\nobreakdash-K}: $\mathbf{K}_{\ell}=\sum_{\Lambda}\omega_{\Lambda}|\mathbf{J}_{\Lambda}|\mathbf{B}_{\Lambda}^\mathrm{T}\mathbf{D}_{\ell}\mathbf{B}_{\Lambda}$.
\item Accumulate nodal forces as
\newline
\textbf{Kernel LMC\nobreakdash-F}: $\mathbf{f}_{\ell}=\mathbf{f}_{\ell}+\mathbf{K}_{\ell}\mathbf{u}_{\ell}$.
\item Write: $\mathbf{f}_{\ell}$ into $\mathbf{f}$.
\end{enumerate}
\end{enumerate}
\item Synchronize $\mathbf{f}$.
\end{enumerate}
\end{algorithm}

%HJ 3. Refusal of the call ---------------------------------------------------
%      the hero experiences some hesitation to answer the call
%
% Transition paragraph draft by Luna Gemini Pro, 2024-09-27, TODO check
Traditional finite element analysis methods often grapple with limitations in handling large-scale problems due to computational bottlenecks associated with matrix storage and evaluation. This necessitates exploring alternative approaches that can overcome these limitations while maintaining computational efficiency.
%TODO Beef this up.

%HJ 4. Meeting with the mentor -----------------------------------------------
%      the hero gains the supplies, knowledge, and confidence needed to 
%      commence the adventure
%
To address this challenge, the subsequent section describes a new approach: representing the element response using tensor trains. Evaluation of these tensor train computational kernels is equivalent to a recent performant low-order matrix-free partial assembly approach for GPU computation.\cite{ref:settgast2023performant}

\section{Development}
\label{sec:Development}

% Intro paragraph draft by Luna Gemini Pro, 2024-09-27, edited by author
This section describes how to derive an efficient tensor train representation of an element's elastic response, starting from the tensor form of the element-by-element method in section~\ref{subsec:EBE-EMS}. A traditional element-by-element local matrix construction method is then recast in terms of tensor algebra. Finally, the modified Einstein summation notation of section~\ref{subsec:einstein-notation} is introduced to simplify tensor expressions. These steps highlight the mathematical basis for the tensor train element response evaluations described in the Methods section~\ref{sec:Methods}.

\subsection{Finite element operations in functional form}
% by the author

Neither the global system matrix nor element stiffness matrices~($\mathbf{K}_{\ell}$) 
need to be constructed when
equation~\ref{eq:EBE-force-assembly} is used to assemble the global response from element responses.
 Instead, stiffness can be applied as an elemental operator~($K_{\ell}$),
\begin{eqnarray}
f_{\ell}(u) = K_{\ell}(u_{\ell})\text{,}\label{eq:functional-elem-response}
\end{eqnarray}
and evaluated as a volume integral of several composed functions,
\begin{eqnarray}
f_{\ell}(u) = \int_{\Omega_{\ell}}B_{\ell\Lambda}^{\top}(D_{\mathcal{P}}(B_{\ell\Lambda}(u_{\ell})))d\Omega_{\ell}\text{.}\label{eq:functional-elem-response-nested}
\end{eqnarray}
Each of the functions composed in equation~\ref{eq:functional-elem-response-nested}
is a tensor analogue of a traditional finite element linear algebra expression.

A strain-displacement~($\mathbf{B}$) matrix can be interpreted as an operator~($B$)
that interpolates the engineering strain~($\vec{\varepsilon}_{\ell\Lambda}$)
from nodal displacements~($u_{\ell}$) at an integration point ($\Lambda$),
\begin{eqnarray}
\vec{\varepsilon}_{\ell\Lambda} = B_{\ell\Lambda}(u_{\ell})\text{,}\label{eq:small-strain}
\end{eqnarray}
As an operator, a material Voigt ($\mathbf{D}$)~matrix calculates the Cauchy stresses from a Voigt vector representation of engineering strains,
\begin{eqnarray}
\vec{\sigma}_{\ell\Lambda} = D_{\mathcal{P}}(\vec{\varepsilon}_{\ell\Lambda})\text{.}
\end{eqnarray}
Finally, the nodal force response equivalent to equation~\ref{eq:EMS-elem-response}
can be evaluated by volume integration of the stresses projected back onto
 the element's shape gradient basis at the nodes by the functional inverse (transpose) of the strain-displacement operator~($B_{\ell\Lambda}^{\top}$),
\begin{eqnarray}
f_{\ell} = \int_{\Omega_{\ell}}B_{\ell\Lambda}^{\top}(\sigma_{\ell\Lambda})d\Omega_{\ell}\text{.}
\end{eqnarray}
These traditional finite element operators are expressions in the linear
 algebra context of the mechanics of materials.
 Equivalent tensor operators can be derived in the more general framework
 of continuum mechanics.
 
\subsection{Finite element equivalent tensor functions}
% by the author

Rather than constructing and applying the traditional sparse ($\mathbf{B}$)~matrix, the shape function gradient operator~($G_{\ell\Lambda}$)
 can be used to evaluate a small strain tensor~($H_{\ell\Lambda}$) at an integration point,
\begin{eqnarray}
H_{\ell\Lambda} = G_{\ell\Lambda}(u_{\ell})\text{.}\label{eq:small-strain-1}
\end{eqnarray}
This computation is more compact than computing engineering strain as a Voigt vector by
equation~\ref{eq:voigt-strain-from-displacement}.
 As an operator, a fourth-order material constitutive tensor~($C_{\mathcal{P}}$)
 calculates the Cauchy stress tensor from a strain tensor.
 The symmetries of an elastic response tensor allow the Cauchy stress at the 
 integration point to be calculated from either the asymmetric small strain tensor~($H_{\ell\Lambda}$)
 or the symmetric true strain tensor~($\epsilon_{\ell\Lambda}$),
\begin{eqnarray}
\sigma_{\ell\Lambda} = C_{\mathcal{P}}(H_{\ell\Lambda})=C_{\mathcal{P}}(\epsilon_{\ell\Lambda})\text{,}
\end{eqnarray}
equivalent to equation~\ref{eq:stress-from-engineering-strain-global}.
 Finally, the nodal force response equivalent to equation~\ref{eq:EMS-elem-response}
can be computed by volume integration of the stresses projected back onto
 the element's shape gradient basis by the functional inverse (transpose) of the shape gradient
 operator~($G_{\ell\Lambda}^{\top}$),

\begin{eqnarray}
f_{\ell} = \int_{\Omega_{\ell}}G_{\ell\Lambda}^{\top}(\sigma_{\ell\Lambda})d\Omega_{\ell}\text{.}
\end{eqnarray}

Thus, the symmetries of a fourth-order elastic response constitutive tensor~($C$)
 allow element stiffness to be computed more compactly than 
 equation~\ref{eq:functional-elem-response-nested} as,
\begin{eqnarray}
f_{\ell}(u) = \int_{\Omega_{\ell}}G_{\ell\Lambda}^{\top}(C_{\mathcal{P}}(G_{\ell\Lambda}(u_{\ell})))d\Omega_{\ell}\text{,}\label{eq:functional-tensor-elem-response}
\end{eqnarray}
when assembling the global response by equation~\ref{eq:EBE-force-assembly}.
 Each of the functions composed in equation~\ref{eq:functional-tensor-elem-response}
 is an equivalent tensor formulation of a traditional finite element linear
 algebra expression.

The finite element method, as a means to solve Navier's elasticity equations,
 is traditionally presented in the context of the mechanics of materials.
 Mechanics of materials is a first-order approximation of elastic behavior
 describing the linear behavior of many engineering materials under
 small strain conditions.
 In this context, non-scalar quantities are represented as vectors such
 as the Voigt vector representations of engineering strains and Cauchy stresses.
 Matrices apply linear operations to these vectors.

% Transition paragraph draft by Luna Gemini Pro, 2024-09-27
While both sparse matrix and traditional EBE methods have been successfully applied, expressing finite element operations in a tensor algebra framework offers valuable insights into potential optimizations. To facilitate this transition, the following subsection introduces a modified Einstein implied summation notation.

\subsection{Modified Einstein implied summation notation}
\label{subsec:einstein-notation}
% by the author

Continuum mechanics generalizes the mechanics of materials and is typically
 presented as tensor operations using the Einstein implied summation convention.\cite{ref:einstein-1922-general}
 Here, repeated italic roman letter subscripts in a term imply summation.
 Lowercase (LC) roman subscripts represent
 the usual 3D~($d_{\mathrm{m}}=3$) continuum spatial basis,
 while uppercase (UC) subscripts represent
 other, non-spatial bases, such as an element's shape function basis.
 Note that, in the context of Einstein notation, these subscripts are arbitrary
 and generally divorced from any meaning assigned in the Nomenclature section.
 
 Script subscripts (e.g., $X_{\mathcal{A}}$, $E_{\ell}$) and 
 Greek letter subscripts (e.g., $K_{\alpha}$) 
 denote identifying indices without implied summation, and upright roman
 letter subscripts (e.g., $Z_{\mathrm{A}}$) are labels. These conventions are summarized in Table~\ref{tab:subscript-notation}.

\begin{table}[ht]
\caption{Notation conventions}
\label{tab:subscript-notation}
\begin{center}
\begin{tabular}{lll}
\toprule
 {Subscript notation} & {Meaning} & {Example} \\
\midrule
 LC italic roman (repeated) & Spatial indices, implied sum & $\sigma_{ij} = C_{ijkl} \epsilon_{kl}$ \\
 UC italic roman (repeated) & non-spatial indices, implied sum & $f_{Aj} = K_{AjBk} u_{Bk}$ \\
 Script & Identifying indices, no sum & $X_{\mathcal{A}}$, $E_{\ell}$ \\
 Greek & Identifying indices, no sum & $K_{\alpha}$ \\
 Upright italic roman & Labels, no sum & $Z_{\mathrm{A}}$ \\
\bottomrule
\end{tabular}
\end{center}
\end{table}

For example, in this form, an element response represented by
equation~\ref{eq:EMS-elem-response} can be expressed as,
\begin{eqnarray}
f_{Bb} = K_{BbaA}u_{Aa}\text{,}\label{eq:EMS-response-tensor}
\end{eqnarray}
with the element identifying index~($\ell$) omitted for clarity.
 The indices repeated in the term on the right hand side of
 equation~\ref{eq:EMS-response-tensor} imply the summation,
\begin{eqnarray}
f_{Bb} = \sum_{A=1}^{n_{\mathrm{n}}}\sum_{a=1}^{d_{\mathrm{m}}}K_{BbaA}u_{Aa}\text{.}\label{eq:explicit-multi-index-sum}
\end{eqnarray}
The tensor multi-indices
($Bb$ and $aA$) can each be transformed into corresponding matrix indices
(${R}$ and ${C}$), recasting the tensor summations of
equation~\ref{eq:explicit-multi-index-sum} into the standard linear algebra expression
for a dense matrix-vector multiplication,
\begin{eqnarray}
f_{R} & = & \sum_{{C}=1}^{n_{\mathrm{e}}}K_{{R}{C}}u_{C}\text{, with} \\
{R} & = & d_{\mathrm{m}}(B-1)+b\text{,} \nonumber \\
{C} & = & d_{\mathrm{m}}(A-1)+a\text{\text{, and}} \nonumber \\
n_{\mathrm{e}} & = & d_{\mathrm{m}}n_{\mathrm{n}}\text{.} \nonumber
\end{eqnarray}
This notation helps recast traditional finite element method linear algebra
 algorithms into equivalent tensor forms presented in the next section.

\section{Methods}
\label{sec:Methods}
%HJ 5. Crossing the first threshold ------------------------------------------
%      the hero commits wholeheartedly to the adventure

% Intro paragrpah draft by Luna Gemini Pro, 2024-09-27
This section details how to efficiently evaluate the physics response of a tetrahedral finite element model using tensor train analysis. This approach computes the global physics response as a sum of element responses, without assembling the global stiffness matrix, as described in section~\ref{subsec:EBE-EMS}. Each element response is evaluated as a tensor train composed of a sequence of small, dense tensor operations. First, an element's response is computed in a coordinate system that simplifies evaluation. Then, if necessary, the response is transformed to global coordinates. Finally, the contributions of all elements are summed to compute the global response.

% Transition paragraph draft by Luna Gemini Pro, 2024-09-27
Before delving into specific element response computations, it's worthwhile to consider simplifying assumptions that can significantly enhance performance without substantial loss of generality.

\subsection{Simplifying assumptions}

Without loss of too much generality, and remaining relevant to many scientific and engineering applications, two limitations improve tetrahedral element response kernel performance by reducing the amount of main memory access required during element response computation.
\begin{enumerate}
\item Models are composed of one or more homogeneous regions (partitions).
\item Quadratic and cubic element side nodes are at their natural co-linear locations between vertex nodes.
\end{enumerate}

When a model is parallelized by domain decomposition, the first limitation
 reduces the material property data that needs to be read from main memory,
 rather than from cache, during element response kernel computation.
 The set of material properties~($c \in \mathcal{C}_{p}$)
 for each homogeneous partition in a domain ($\mathcal{P}_{p} \in \Omega$) can
 be read once and then reused for all of the elements in that partition,
\begin{equation}
\forall\ell\in\mathcal{P}_{p}:\mathcal{C}_{\ell\Lambda} = \mathcal{C}_{p}\text{.}
\end{equation}

The second limitation reduces the amount of geometric data needed for each
 element response evaluation.
 When a triangular or tetrahedral element's side nodes are at their natural
 co-linear locations, the element's distortion from its reference configuration
 is only a function of the positions of its vertex nodes.
 The element's geometric Jacobian can be computed from the vertex nodal
 coordinates~($\mathbf{x}_{\ell}^{\mathrm{P1}}$) and the element's 
 first order~(P1) shape function derivatives~($\dot{\mathbf{N}}_{\Lambda}^{\mathrm{P1}}$)
 at the integration points.
 In addition, first-order triangular and tetrahedral element geometric distortion
 is constant throughout the element.
 This implies that any triangular or tetrahedral element's geometric
 Jacobian~($\mathbf{J}_{\ell}$), and its inverse, are the same at every 
 integration point~($\Lambda$) within the element,
\begin{equation}
\forall\Lambda\in\ell:\mathbf{J}_{\ell} = \mathbf{J}_{\ell\Lambda}^{\mathrm{P1}} = \dot{\mathbf{N}}_{\Lambda}^{\mathrm{P1}} \mathbf{x}_{\ell}^{\mathrm{P1}}\text{.}
\end{equation}
Thus, only one inverted Jacobian is needed to characterize the geometry
 of each triangular or tetrahedral element when its side nodes (if present) are at their
 natural positions.

The computational cost, memory bandwidth, and memory latency of element
 operations can be reduced by initializing an array of concatenated transposed inverted element geometric 
 Jacobians~($\mathbf{J}_{\ell}^{-\top}$) and Jacobian determinants~($|\mathbf{J}_{\ell}|$).
 These can be stored for uniform read access during element response
 evaluations rather than the non-uniform node coordinate access needed to compute the Jacobian for each element.

% Transition paragraph draft by Luna Gemini Pro, 2024-09-27, edited by the author
%TODO FIXME remove "we"
Having established a tensor notation and some simplifying assumptions, we can now re-examine the computation of element responses. The following subsection details how this is achieved for elastic isotropic materials or materials aligned with the global coordinates.

\subsection{Elastic element response in global coordinates}
\label{subsec:EBE-tensor-train-iso}

Typical finite element implementations use a strain-displacement~($\mathbf{B}$)
matrix to obtain engineering strain~($\vec{\varepsilon}$) at an integration point
within a solid elastic element, as shown in equation~\ref{eq:voigt-strain-from-displacement}.
However, an equivalent asymmetric small strain tensor~($H_{\ell\Lambda}$)
can be calculated from the element's nodal displacements~($u_{\ell}$)
with a more compact tensor operation,
\begin{eqnarray}
H_{lk} & = & G_{lB}u_{Bk}\text{.}\label{eq:tensor-strain-from-displacement}\\
\left[\begin{array}{ccc}
\frac{\partial u_{x}}{\partial x} & \frac{\partial u_{y}}{\partial x} & \frac{\partial u_{z}}{\partial x}\\
\frac{\partial u_{x}}{\partial y} & \frac{\partial u_{y}}{\partial y} & \frac{\partial u_{z}}{\partial y}\\
\frac{\partial u_{x}}{\partial z} & \frac{\partial u_{y}}{\partial z} & \frac{\partial u_{z}}{\partial z}
\end{array}\right] & = & \left[\begin{array}{ccc}
\frac{\partial N_{1}}{\partial x} &  & \frac{\partial N_{n_{n}}}{\partial x}\\
\frac{\partial N_{1}}{\partial y} & \cdots & \frac{\partial N_{n_{n}}}{\partial y}\\
\frac{\partial N_{1}}{\partial z} &  & \frac{\partial Nn_{n}}{\partial z}
\end{array}\right]\left[\begin{array}{ccc}
u_{\ell1x} & u_{\ell1y} & u_{\ell1z}\\
 & \vdots\\
u_{\ell n_{n}z} & u_{\ell n_{n}y} & u_{\ell n_{n}z}
\end{array}\right]\text{.}\nonumber 
\end{eqnarray}

This second-order small strain tensor is the material Jacobian at the integration
 point in global coordinates.
 Compare equations~\ref{eq:small-strain} to the geometric Jacobian calculation in equation~\ref{eq:geometric-jacobian}.
 This small strain tensor is directly related to the continuum mechanics
 material deformation tensor~($F$), which can be obtained by adding a conformal identity tensor~($\delta$), 
\begin{eqnarray}
F_{lk} = H_{lk}+\delta_{lk}\text{.}
\end{eqnarray}
If needed, the symmetric true strain tensor~($\epsilon$)
 and engineering strain Voigt vector~($\vec{\varepsilon}$) at the integration point can also
 be calculated from the asymmetric small strain tensor as,
\begin{eqnarray}
\epsilon_{lk} & = & \tfrac{1}{2}(H^{\top}_{kl}+H_{lk})_{lk}\text{, and}\label{eq:true-strain-from-unsymm-strain}\\
\vec{\varepsilon} & = & \left[\begin{array}{cccccc}
H_{11}, & H_{22}, & H_{33}, & H_{23}+H_{32}, & H_{13}+H_{31}, & H_{12}+H_{21}\end{array}\right]^{\top}\text{.}\label{eq:engineering-strain-from unsymm-strain}
\end{eqnarray}

The symmetries of a fourth-order elastic material response (constitutive)
 tensor~($C$) from continuum mechanics allow the Cauchy stress tensor~($\sigma$) to be calculated
 from either the asymmetric small strain or the symmetric true strain as,
\begin{eqnarray}
\sigma_{ij} = C_{ijkl} H_{lk} = C_{ijkl} \epsilon_{lk}\text{.}\label{eq:tensor-stress-from-strain}
\end{eqnarray}
These symmetries, and the sparse structure of constitutive tensors for many
 common engineering materials, can be exploited to improve performance,
 as detailed in the next section.% TODO explain c_{unq}
 
The nodal force response can be computed by numerically integrating
 the stresses within the element,
\begin{eqnarray}
f_{Aj} = \sum_{\Lambda} \omega_{\Lambda} |J_{\ell}| G_{Ai}^{\top} \sigma_{ij}\label{eq:tensor-forces-from-stress-integration}
\end{eqnarray}
Composing the functional form of equation~\ref{eq:functional-tensor-elem-response} by substituting equation~\ref{eq:tensor-strain-from-displacement} into~\ref{eq:tensor-stress-from-strain}, and the result
 into~\ref{eq:tensor-forces-from-stress-integration}, yields the element response in implied summation notation as,
\begin{eqnarray}
f_{Aj} = \sum_{\Lambda} \omega_{\Lambda} |J_{\ell}| G_{Ai}^{\top} C_{ijkl} G_{lB} u_{Bk}\text{.}\label{eq:tensor-elem-response-elastic}
\end{eqnarray}
Embedded within equation~\ref{eq:tensor-elem-response-elastic} is a tensor expression of element stiffness, 
\begin{eqnarray}
K_{AabB} = \sum_{\Lambda}\omega_{\Lambda} |J_{\ell}| G_{Ai} C_{iabj} G_{jB}\text{,}\label{eq:tensor-elem-stiff}
\end{eqnarray}
equivalent to the linear algebra element stiffness matrix computation of
 equation~\ref{eq:la-elem-stiff}, which can be seen by comparing 
 equation~\ref{eq:tensor-elem-response-elastic} to equation~\ref{eq:EMS-response-tensor}.

Evaluating the element response tensor expression, equation~\ref{eq:tensor-elem-response-elastic}, from right to left, or equation~\ref{eq:functional-tensor-elem-response} from inside out, suggests Algorithm~\ref{alg:iso-elastic-system-response} for computing the system response of a linear-elastic model.
\begin{algorithm}[ht]%WAS hbt!
%\caption{Latency-hiding EBE tensor train response evaluation in global coordinates}
\caption{Global coordinates elastic response tensor kernels}
\label{alg:iso-elastic-system-response}
\begin{enumerate}
\item Initialize:
Transposed inverted Jacobians~(${J}_{\ell}^{-\top}$) and determinants~($|{J}_{\ell}|$).
\item For each homogeneous partition~($\mathcal{P}$):
\begin{enumerate}
\item Read: $\dot{N}$, $\omega$, $c_{\mathrm{unq}}(C_{\mathcal{P}})$,
${u}_{0}$, ${{J}}_{0}^{-\top}$, $|{J}_{0}|$.
\item For each halo element~($\ell \in \mathcal{E}_{\mathrm{halo}}$):
\begin{enumerate}
\item Read: $f_{Aj}$ from ${f}$.
\item For each integration point~($\Lambda$):
\begin{enumerate}
\item Compute the shape function gradient and small strain in\break global coordinates,
\newline
Kernel~\textbf{GLB\nobreakdash-GH}: $G_{Ai}^{\top} = \dot{N}_{\Lambda Aj}^{\top} {J}_{ji}^{-\top}\text{ and }H_{ia} = G_{iA} u_{Aa}$.
\item Calculate weighted differential volume: $\iota = |{J}_{\ell}| \omega_{\Lambda}$.
\item If this is the last integration point and not the last element:\break prefetch the next ${u}_{\ell+1}$, ${J}_{\ell+1}^{-\top}, |{J}_{\ell+1}|$.
\item Compute weighted differential Cauchy stress,
\newline
Kernel~\textbf{GLB\nobreakdash-S}: $\sigma_{ij}^{\iota} = \iota C(c_{unq})_{ijkl} H_{lk}$.
\item Accumulate nodal forces.
\newline
Kernel~\textbf{GLB\nobreakdash-F}: $f_{Aj} = f_{Aj} + G_{Ai}^{\top}\sigma_{ij}^{\iota}$.
\end{enumerate}
\item Write: $f_{Aj}$ into ${f}$.
\end{enumerate}
\item Synchronize halo node forces ($f \in \mathcal{F}_{\mathrm{halo}}$).
\item Repeat~(b) for interior elements ($\ell \notin \mathcal{E}_{\mathrm{halo}}$).
\end{enumerate}
\end{enumerate}
\end{algorithm}

Computational complexity

Kernel GLB-GH: $O(n_n \times d_m^2)$ operations %TODO check

Kernel GLB-S: $O(d_m^3)$ operations (depends on material symmetry)%TODO check

Kernel GLB-F: $O(n_n \times d_m^2$) operations %TODO check

Per element: $O(n_\Lambda \times n_n \times d_m^2)$ total %TODO check

Memory access patterns

Sequential reads: $J^{-\top}$

Random access reads: $u$

Random access writes: $f$ (with atomic adds for shared nodes)

Constant data: $\dot{N}$, $\omega$, $C$


This algorithm hides some of the time needed to access the input by requesting
 data for the next element before completing the computations needed to
 evaluate the current element.
 Isotropic material response is insensitive to rotations, but it may be
 more convenient to evaluate anisotropic material response in material coordinates.

% Transition paragraph draft by Luna Gemini Pro, 2024-09-27
While formulating the elastic element response in global coordinates is convenient for isotropic materials, anisotropic materials often simplify in their material coordinate system. Therefore, the following subsection describes how to efficiently compute the element response directly in material coordinates. [TODO FIXME]

\subsection{Elastic element response in material coordinates}
\label{subsec:EBE-tensor-train-aniso}

For anisotropic materials, a rotation tensor~($R$) transforms a fourth-order constitutive operator in material coordinates,~$\widetilde{C}(\widetilde{\epsilon})$ to global coordinates,~$C(\epsilon)$, as,
% \begin{eqnarray}
% C_{nboa} & = & R_{ni}^{\top} R_{bj}^{\top} \widetilde{C}_{ijkl} R_{ko} R_{la}\text{.}\label{eq:rotate-material-tensor}
% \end{eqnarray}
\begin{eqnarray}
\label{eq:rotate-material-tensor}
C_{ijkl} = R_{ia}^{\top}R_{jb}^{\top}\widetilde{C}_{abcd}R_{dk}R_{cl}\text{.}
\end{eqnarray}
Substituting equation~\ref{eq:rotate-material-tensor} into equation~\ref{eq:tensor-elem-response-elastic} yields,
% \begin{eqnarray}
% f_{Bb} & = & \sum_{\Lambda}\omega_{\Lambda}|J|\dot{N}_{Bh}^{\top}J_{hn}^{-\top}R_{ni}^{\top}R_{bj}^{\top}\widetilde{C}_{ijkl}R_{la}R_{ko}J_{om}^{-1}\dot{N}_{mA}u_{Aa}\text{.}\label{eq:elem-response-elastic-ortho}
% \end{eqnarray}
\begin{eqnarray}\label{eq:elem-response-elastic-ortho}
f_{Aj} & = & \sum_{\Lambda}\omega_{\Lambda}|J| \dot{N}_{Ae}^{\top} J_{ei}^{-\top} R_{ia}^{\top} R_{jb}^{\top} \widetilde{C}_{abcd} R_{dk} R_{cl} J_{lf}^{-1} \dot{N}_{fB} u_{Bk}\text{, with}\\
%FIXME TODO definition is: \widetilde{C}_{abcd} = R_{ai} R_{bj} R_{ck} R_{dl} \C_{ijkl}. Above is transposed.
%TODO Explain why, is material orientation defined as the inverse?
{G}_{lB} & = & {J}_{lf}^{-1}\dot{N}_{fB}\text{, and}\\
{G}^{\top}_{Ai} & = & \dot{N}_{Ae}^{\top}{J}_{ei}^{-\top}\text{.}
\end{eqnarray}
Equation~\ref{eq:elem-response-elastic-ortho} suggests an anisotropic response computational kernel,
\begin{eqnarray}
\label{eq:geo-distort-rotated}
f_{Aj} & = & %\sum_{\Lambda} \omega_{\Lambda} |J| R_{jb}^{\top} \widetilde{G}_{Aa}^{\top} \widetilde{C}_{abcd}R_{dk} \widetilde{G}_{cB} u_{Bk}\text{, with}\label{eq:elem-response-elastic-ortho-compact}\\
\sum_{\Lambda} \omega_{\Lambda} |J| \widetilde{G}_{Aa}^{\top} \widetilde{C}_{abcd} \breve{\epsilon}_{ck} R_{kd}^{\top} R_{bj}\text{, with}\label{eq:elem-response-elastic-ortho-compact}\\
%\sum_{\Lambda}\omega_{\Lambda}|J|R_{jb}^{\top}\widetilde{G}_{Aa}^{\top}\widetilde{C}_{abcd}R_{dk}\widetilde{G}_{cB}u_{Bk}\text{, with}\label{eq:elem-response-elastic-ortho-compact}\\
\widetilde{J}_{cf}^{-1} & = & R_{cl} J_{lf}^{-1}\text{,}\label{eq:jacobian-rotated}\\
\widetilde{G}_{cB} & = & \widetilde{J}_{cf}^{-1} \dot{N}_{fB}\text{, and}\\
\breve{\epsilon}_{ck} & = & \widetilde{G}_{cB} u_{Bk}\text{.}\label{eq:aniso-temp-variable}
\end{eqnarray}
Algorithm~\ref{alg:Anisotropic-linear-elastic} implements
Equations~\ref{eq:elem-response-elastic-ortho-compact}\nobreakdash-\ref{eq:aniso-temp-variable}.

\begin{algorithm}[ht]%WAS hbt!
\caption{Material coordinates elastic response tensor kernels}
\label{alg:Anisotropic-linear-elastic}
\begin{enumerate}
\item Initialize:
Rotated inverted Jacobians~($\widetilde{J}_{{\ell}0}^{-1}$) and determinants~($|{J}_{\ell}|$).
\item For each homogeneous partition:
\begin{enumerate}
\item Read: $\dot{N}$, $\omega$, $c_\text{unq}(\widetilde{C}_{\mathcal{P}}), R$,
$u_{{\ell}0}$, $|J_{{\ell}0}|$, $\widetilde{J}_{{\ell}0}^{-1}$.
%\item For each tetrahedral element:
\item For each halo element~($\ell \in \mathcal{E}_{\mathrm{halo}}$):
\begin{enumerate}
\item Read: $f_{\ell}$.
\item For each integration point:
\begin{enumerate}
\item Read: $\dot{N}_{fB} = \dot{N_{\Lambda}}, \omega_{\Lambda}$.
\item Calculate weighted differential volume: $\iota = |J_{\ell}| \omega_{\Lambda}$.
\item Compute shape function gradient and asymmetric strain\break in material coordinates,

Kernel \textbf{MTR\nobreakdash-GH}: $\widetilde{G}_{cB} = \widetilde{J}_{cf}^{-1} \dot{N}_{fB}$, $\breve{\epsilon}_{ck} = \widetilde{G}_{cB} u_{Bk}$, and $\widetilde{H}_{cd} = \breve{\epsilon}_{ck} R_{kd}^{\top}$.
\item If this is the last integration point and not the last element:\break prefetch the next $u_{\ell+1}$, $J_{\ell+1}^{-1}, |J_{\ell+1}|$.
\item Compute weighted stress in material coordinates:

Kernel \textbf{MTR\nobreakdash-S}: $\widetilde{\sigma}_{ab}^{\omega} = \iota \widetilde{C}(c_\text{unq})_{abcd} \widetilde{H}^{\top}_{dc}$.
\item Accumulate nodal forces,

Kernel \textbf{MTR\nobreakdash-F}: $b_{aj} = \widetilde{\sigma}_{ab}^{\omega} R_{bj}$ and $f_{Aj} = f_{Aj} + \widetilde{G}_{Aa}^{\top} b_{aj}$.
\end{enumerate}
\item Write: $f_{e}$.
\end{enumerate}
\item Synchronize halo node forces ($f \in \mathcal{F}_{\mathrm{halo}}$).
\item Repeat~(b) for interior elements ($\ell \notin \mathcal{E}_{\mathrm{halo}}$).
\end{enumerate}
\end{enumerate}
\end{algorithm}

Rotated Jacobians~($\widetilde{J}^{-1}$)
 can be computed by equation~\ref{eq:jacobian-rotated} and stored for each element during initialization, before
 evaluating this kernel.
 Note that the rotations remaining in equation~\ref{eq:elem-response-elastic-ortho-compact} have been transposed from their representations in equation~\ref{eq:elem-response-elastic-ortho}.
 An asymmetric small strain tensor~($\widetilde{H}$) in material coordinates can be computed as,
\begin{eqnarray}
\widetilde{H}_{cd} = \widetilde{G}_{cB} u_{Bk} R_{kd}^{\top}\text{.}
\end{eqnarray}
The Cauchy stress in material coordinates~($\widetilde{\sigma}_{ab}$) 
 comes directly from a material response evaluation,
\begin{eqnarray}
\widetilde{\sigma}_{ab} = \widetilde{C}_{abcd} \widetilde{H}^{\top}_{dc}\text{.}
\end{eqnarray}


\section{Implementation details}
\label{sec:Implementation}

\subsection{Data structures}
\label{subsec:data-structures}


\subsection{Vectorization}
\label{subsec:vectorization}

\subsection{Compiler optimization}
\label{subsec:compiler-optimization}

\subsection{Parallelization}
\label{subsec:parallelization}

% Transition draft by Luna Gemini Pro, 2024-09-15
Having established an iterative solution strategy and efficient evaluation of the global physical response, the next section presents concrete examples to demonstrate the capabilities and performance of a parallel EBERTTA finite element analysis~(FEA) framework.

\section{Examples}
\label{sec:Examples}
%HJ 7. Approach to the inmost cave -------------------------------------------
%      the hero nears the center of the story and the special world


%TODO Verification: NAFEMS Finite Element Benchmarks for MDG Code Verification?
%TODO Patch tests

% Intro draft by Luna Gemini Pro, 2024-09-15
These examples highlight the effectiveness of the EBERTTA method in solving diverse problems with varying complexities and scales. The demonstrated scalability and performance gains make this FEA framework well-suited for tackling large-scale scientific and engineering simulation efforts. 
Source code is available in the Femera mini-app~v0.1.~\cite{ref:femera-v0.1-source} 
% TODO rewrite transition?

\subsection{General linear elastic (triclinic)}

%TODO material symmetries and sparseness

Symmetries of an elastic material constitutive tensor~($\widetilde{c}$) reduce the number of unique values needed to characterize the material to 21. This allows recasting the stress-strain relationship into linear algebra using a symmetric $6\times6$ Voigt matrix.
%to transform strains to stresses
%and arranging the vectors
%as shown in equation~\ref{eq:mtrl-d-matrix} 
For evaluation in material coordinates, kernel CMV~S,
 
\begin{eqnarray}
\widetilde{\vec{\mathbf{\sigma}}} & = & \widetilde{\mathbf{D}}\widetilde{\vec{\mathbf{\varepsilon}}}\text{,}\label{eq:mtrl-d-matrix}{}\\
\left[\begin{array}{c}
\widetilde{\sigma}_{x}\\
\widetilde{\sigma}_{y}\\
\widetilde{\sigma}_{z}\\
\widetilde{\sigma}_{yz}\\
\widetilde{\sigma}_{xz}\\
\widetilde{\sigma}_{xy}
\end{array}\right] & = & \left[\begin{array}{cccccc}
%
\widetilde{c}_{11} & & & & &\\
\widetilde{c}_{21} & \widetilde{c}_{22} & & \multicolumn{2}{c}{\text{symm.}} &\\
\widetilde{c}_{31} & \widetilde{c}_{32} & \widetilde{c}_{33} & & &\\
\widetilde{c}_{41} & \widetilde{c}_{42} & \widetilde{c}_{43} & \widetilde{c}_{44} & &\\
\widetilde{c}_{51} & \widetilde{c}_{52} & \widetilde{c}_{53} & \widetilde{c}_{54} & \widetilde{c}_{55} &\\
\widetilde{c}_{61} & \widetilde{c}_{62} & \widetilde{c}_{63} & \widetilde{c}_{64} & \widetilde{c}_{64} &  \widetilde{c}_{66}
%
%  \widetilde{c}_{11} & \widetilde{c}_{12} & \widetilde{c}_{13} & \widetilde{c}_{14} & \widetilde{c}_{15} & \widetilde{c}_{16}\\
% & \widetilde{c}_{22} & \widetilde{c}_{22} & \widetilde{c}_{23} & \widetilde{c}_{24} & \widetilde{c}_{25} \\
% & & \widetilde{c}_{33} & \widetilde{c}_{34} & \widetilde{c}_{34} & \widetilde{c}_{36} \\
% & & & \widetilde{c}_{44} & \widetilde{c}_{45} & \widetilde{c}_{46} \\
% & \multicolumn{2}{c}{\text{symm.}} & &\widetilde{c}_{55} & \widetilde{c}_{56} \\
% & & & & &\widetilde{c}_{66}
%
\end{array}\right]\left[\begin{array}{c}
\widetilde{\varepsilon}_{x}\\
\widetilde{\varepsilon}_{y}\\
\widetilde{\varepsilon}_{z}\\
\widetilde{\gamma}_{yz}\\
\widetilde{\gamma}_{xz}\\
\widetilde{\gamma}_{xy}\nonumber\label{eq:stress-from-engineering-strain-material}
\end{array}\right]\text{.}
\end{eqnarray}
is a dense symmetric matrix-vector product, with
engineering strain calculated by equation~\ref{eq:engineering-strain-from unsymm-strain}.
However, many structural materials of engineering interest are orthotropic (including transverse isotropic, cubic, and isotropic),
and can be characterized by at most nine independent material properties~($\widetilde{c}$)
when the material coordinates are aligned with orthogonal symmetry planes. 

\subsection{Orthotropic linear elastic}
% Luna Gemini Pro 2024-09-12
% Orthotropic materials exhibit distinct material properties along three mutually perpendicular planes of symmetry, requiring a maximum of nine independent material constants for characterization when material coordinates align with these planes. Cubic materials, a subset of orthotropic materials, possess equal properties along all three axes of symmetry, simplifying their description to just three independent constants. Transversely isotropic materials exhibit rotational symmetry about one axis, reducing the necessary material parameters to five. Isotropic materials, representing the most symmetric case, possess uniform properties in all directions, necessitating only two independent constants for characterization. Notably, the constitutive tensor structure for orthotropic, cubic, transversely isotropic, and isotropic materials remains consistent, allowing isotropic materials to be treated as a special case of orthotropic response. This shared structure facilitates efficient computation of elastic response across these material classes.

% Author draft
Orthotropic (9~unique properties) materials % have the same sparse constitutive tensor structure as isotropic tensors and 
include transversely isotropic (5~properties) and cubic (3~properties) material models. Isotropic materials (2~properties) can also be considered as a special case of orthotropic response because the material response tensor has the same sparse structure. This sparse structure can be represented as a Voigt matrix, %TODO reqrite this.
%
\begin{eqnarray}
\widetilde{\mathbf{D}}_{\mathrm{ort}} & = & \left[\begin{array}{cccccc}
%
\widetilde{c}_{11} & & & & &\\
\widetilde{c}_{21} & \widetilde{c}_{22} & & \multicolumn{2}{c}{\text{symm.}} &\\
\widetilde{c}_{31} & \widetilde{c}_{32} & \widetilde{c}_{33} & & &\\
0 & 0 & 0 & \widetilde{c}_{44} & &\\
0 & 0 & 0 & 0 & \widetilde{c}_{55} &\\
0 & 0 & 0 & 0 & 0 &  \widetilde{c}_{66}
\end{array}\right]\text{.}
\end{eqnarray}

The elastic response of orthotropic materials in material coordinates can be computed more efficiently using equations~\ref{eq:material-response-ortho},
\begin{eqnarray}
\widetilde{\sigma}_{x} = \widetilde{\sigma}_{11} & = & \widetilde{c}_{11}\widetilde{H}_{11} + \widetilde{c}_{12}\widetilde{H}_{22} + \widetilde{c}_{13}\widetilde{H}_{33}\text{,}\label{eq:material-response-ortho}\\
\widetilde{\sigma}_{y} = \widetilde{\sigma}_{22} & = & \widetilde{c}_{12}\widetilde{H}_{11} + \widetilde{c}_{22}\widetilde{H}_{22} + \widetilde{c}_{23}\widetilde{H}_{33}\text{,}\nonumber \\
\widetilde{\sigma}_{z} = \widetilde{\sigma}_{33} & = & \widetilde{c}_{13}\widetilde{H}_{11} + \widetilde{c}_{23}\widetilde{H}_{22} + \widetilde{c}_{33}\widetilde{H}_{33}\text{,}\nonumber \\
\widetilde{\sigma}_{yz} = \widetilde{\sigma}_{23} = \widetilde{\sigma}_{32} & = & \widetilde{c}_{44}(\widetilde{H}_{23} + \widetilde{H}_{32})\text{,}\nonumber \\
\widetilde{\sigma}_{xz} = \widetilde{\sigma}_{13} = \widetilde{\sigma}_{31} & = & \widetilde{c}_{55}(\widetilde{H}_{13} + \widetilde{H}_{31})\text{, and}\nonumber \\
\widetilde{\sigma}_{xy} = \widetilde{\sigma}_{12} = \widetilde{\sigma}_{21} & = & \widetilde{c}_{66}(\widetilde{H}_{12} + \widetilde{H}_{21})\text{.}\nonumber
\end{eqnarray}

\subsection{Cubic and isotropic}
Cubic elastic material is characterized by three independent parameters,
\begin{eqnarray}
\widetilde{c}_{1} & = & \widetilde{c}_{11}=\widetilde{c}_{22}=\widetilde{c}_{33}\text{,} \label{eq:cubic-3-parameters} \\
\widetilde{c}_{2} & = & \widetilde{c}_{23}=\widetilde{c}_{13}=\widetilde{c}_{12}=\widetilde{c}_{32}=\widetilde{c}_{31}=\widetilde{c}_{21}\text{, and}\nonumber \\
\widetilde{c}_{3} & = & \widetilde{c}_{44}=\widetilde{c}_{55}=\widetilde{c}_{66}\text{.} \nonumber
\end{eqnarray}
Cubic material response can be evaluated in material coordinates by equations~\ref{eq:material-response-cubic},
\begin{eqnarray}
\widetilde{\sigma}_{11} & = & \widetilde{c}_{1}H_{11} + \widetilde{c}_{2}H_{22} + \widetilde{c}_{2}H_{33}\text{,}\label{eq:material-response-cubic}\\
\widetilde{\sigma}_{22} & = & \widetilde{c}_{2}\widetilde{H}_{11} + \widetilde{c}_{1}\widetilde{H}_{22} + \widetilde{c}_{2}\widetilde{H}_{33}\text{,}\nonumber \\
\widetilde{\sigma}_{33} & = & \widetilde{c}_{2}\widetilde{H}_{11} + \widetilde{c}_{2}\widetilde{H}_{22} + \widetilde{c}_{1}\widetilde{H}_{33}\text{,}\nonumber \\
\widetilde{\sigma}_{23}=\widetilde{\sigma}_{32} & = & \widetilde{c}_{3}\widetilde{H}_{23} + \widetilde{c}_{3}\widetilde{H}_{32}\text{,}\nonumber \\
\widetilde{\sigma}_{13}=\widetilde{\sigma}_{31} & = & \widetilde{c}_{3}\widetilde{H}_{13} + \widetilde{c}_{3}\widetilde{H}_{31}\text{, and}\nonumber \\
\widetilde{\sigma}_{12}=\widetilde{\sigma}_{21} & = & \widetilde{c}_{3}\widetilde{H}_{12} + \widetilde{c}_{3}\widetilde{H}_{21}\text{.} \nonumber
\end{eqnarray}

The symmetries of an isotropic material response tensor allow the Cauchy stress tensor~($\sigma$) to be calculated in any Cartesian coordinate system from either the asymmetric small strain or the symmetric true strain tensor using the two Lamé elastic parameters~($\lambda, \mu$) as,
\begin{eqnarray}
\sigma_{ij} = C_{ijkl}H_{lk}=C_{ijkl}\epsilon_{lk}=\mu\left(H_{lk}+H_{kl}\right)_{ij}+\lambda\delta_{ij}H_{kk}\text{.}\label{eq:material-response-iso-lame}
\end{eqnarray}
Equivalently, for materials having Young's modulus~($E$) and Poisson's ratio~($\nu$), 
 isotropic material response can also be evaluated in global or material coordinates by equations~\ref{eq:material-response-cubic} with,
\begin{eqnarray}
c_{1} & = & \widetilde{c}_{1} = \lambda + 2\mu=\frac{(1-\nu)E}{(1+\nu)(1-2\nu)}\text{,} \label{eq:material-response-iso-E-nu} \\
c_{2} & = & \widetilde{c}_{2} = \lambda=\frac{\nu E}{(1+\nu)(1-2\nu)}\text{, and} \nonumber \\
c_{3} & = & \widetilde{c}_{3} = \mu=\frac{E}{2(1+\nu)}\text{.} \nonumber
\end{eqnarray}

Switching a processor's floating point unit~(FPU) between scalar and vector modes of operation within a computational kernel can take a significant amount of time.\cite{ref:fpu-scalar-vector-switch-time} Because of this,
equation~\ref{eq:material-response-iso-E-nu} is generally more suitable for scalar evaluation,
while a vectorized implementation of equation~\ref{eq:material-response-iso-lame}
can perform better on processors having at least advanced vector extensions~(AVX) single-instruction multiple data~(SIMD) capabilities. 


TODO

\begin{figure}[ht]%WAS h
\begin{center}
\begin{minipage}[c]{0.66\linewidth}
\begin{center}
\includegraphics[width=\linewidth]{fig/femera-26MDOF-6148.eps}
\end{center}
\end{minipage}
\begin{minipage}[c]{0.33\linewidth}
\begin{center}
\includegraphics[width=\linewidth]{fig/LoF-pores.png}
\end{center}
\end{minipage}
\caption{Lack of fusion pore model:
linear elastic isotropic solution time and memory used (Skylake~A)(left),
and visualization of the strain field (right).
}
\label{fig:case-elas-iso-lof-pore}
\end{center}
\end{figure}

\begin{table}[ht]
\caption{Linear elastic kernel performance}
\label{tab:linear-elastic-kernel-performance}
\begin{center}
\begin{tabular}{lSll}
\toprule
\makecell{Material\\model} & \mcone{c}{\makecell{Independent\\properties} } & \makecell{Kernel\\complexity} & \makecell{Typical\\performance} \\
\midrule
 General     & 21 & Dense symmetric MV & Baseline \\
 Orthotropic &  9 & Sparse operations  & $\approx X \times$ faster \\
 Cubic       &  3 & Very sparse        & $\approx Y \times$ faster \\
 Isotropic   &  2 & Minimal            & $\approx Z \times$ faster \\
\bottomrule
\end{tabular}
\end{center}
\end{table}

TODO transition


\subsection{Nonlinear elastic-plastic}
\label{subsec:Nonlinear-Elastic-plastic}

%\subsection{Nonlinear elastic-plastic damage accumulation with cohesive zones}
%\label{subsec:Nonlinear-Elastic-plastic-Damage}
%
%TODO

TODO

\subsection{Linear thermo-elastic}
\label{subsec:Linear-Thermo-elastic}

TODO

\subsection{Linear thermal}
\label{subsec:Linear-Thermal}

TODO

\section{Discussion}
\label{sec:Discussion}
%HJ 8. The ordeal ------------------------------------------------------------
%      the hero faces the greatest challenge yet and experiences death and 
%      rebirth

TODO


\begin{table}[ht]%WAS h
\caption{Tetrahedral finite elements for elastic models}
\label{tab:tet-types}
\begin{center}
\begin{tabular}{rcSSSS}
\toprule
%& {Element} & \mcone{c}{Number} & \mcone{c}{Integration} & \multicolumn{2}{c}{Typical}\\
%& {polynomial} & \mcone{c}{of nodes} & \mcone{c}{points} & \multicolumn{2}{c}{mesh connectivity}\\
& \mcone{c}{\makecell{Element\\polynomial} }  & \mcone{c}{\makecell{Number\\of nodes} } & \mcone{c}{\makecell{Integration\\points} } & \multicolumn{2}{c}{\makecell{Typical\\mesh connectivity} }\\
& type & \mcone{c}{($n_{n}$)} & \mcone{c}{($n_{\Lambda}$)} & \mcone{c}{(elem/dof)} & \mcone{c}{(nnz/dof)}\\
\midrule
Linear & P1 & 4 & 1 & 2.02 & 40.5\\
Quadratic & P2 & 10 & 4 & 0.24 & 85.1\\%TODO CHECK P2 or T2?
Cubic & P3 & 20 & 10 & 0.072 & 142.4\\
\bottomrule
\end{tabular}
\end{center}
\end{table}

TODO

\subsection{Linear elastic isotropic single-node performance}
\label{subsec:Linear-elastic-Isotropic}

TODO

\begin{figure}[ht]%WAS h
\begin{center}
\begin{minipage}[c]{0.48\linewidth}
\begin{center}
\includegraphics[width=\linewidth]{fig/femera-profile-elas-iso-6148b-gcc.eps}
%\caption{
Single-node performance profiles for linear, quadratic, and cubic tetrahedral models
%}
\label{fig:profile-tet10-6148}
\end{center}
\end{minipage}
\begin{minipage}[c]{0.48\linewidth}
\begin{center}
\includegraphics[width=\linewidth]{fig/femera-core-time-6148-iso.eps}
%\caption{
Quadratic tetrahedral model kernel time
%}
\end{center}
\end{minipage}
\caption{Isotropic elastic performance (Skylake~B)}
\end{center}
\end{figure}

Table~\ref{tab:advisor-meshes} details meshes used to characterize elasticity kernel performance at small (5~kdof), medium (500~kdof), and large (50~Mdof) system sizes. The open-source application Gmsh~\cite{ref:gmsh} 4.4.0 generated the meshes used for these analyses, and decomposed the domain into fewer than 2000 partitions with its METIS~\cite{ref:karypis-1997-metis} capability. A simple orthogonal slicing algorithm decomposed larger test models into 2000 or more partitions. 

\begin{table}[ht]%WAS h
\caption{Typical tetrahedral meshes for kernel performance testing}
\label{tab:advisor-meshes}
\begin{center}
\begin{tabular}{rrrrrrrrrr}
\toprule
& \multicolumn{3}{c}{Small} & \multicolumn{3}{c}{Medium} & \multicolumn{3}{c}{Large}\\
Nominal Size & \multicolumn{3}{c}{5 kdof} & \multicolumn{3}{c}{500 kdof} & \multicolumn{3}{c}{50 Mdof}\\
Partitions & \multicolumn{3}{c}{1} & \multicolumn{3}{c}{240 (280)} & \multicolumn{3}{c}{3360 and $n_{cpu}$}\\
Test Models & \multicolumn{3}{c}{240 (280)} & \multicolumn{3}{c}{1} & \multicolumn{3}{c}{1}\\
Iterations & \multicolumn{3}{c}{4167 (3571)} & \multicolumn{3}{c}{10000} & \multicolumn{3}{c}{100}\\
\cmidrule(lr){2-4}\cmidrule(lr){5-7}\cmidrule(lr){8-10}
Element Type & P1 & P2 & P3 & P1 & P2 & P3 & P1 & P2 & P3\\
Elements\\
Unique Nodes & 1959 & 1707 & 1452 & 160273 & 175652 & 166990 & 16555928 & 16314014 & 16072183\\
Model Size & 5877 & 5121 & 4356 & 480819 & 526956 & 500970 & 49667784 & 48942042 & 48216549\\
\cmidrule(lr){2-10}
Total Test Size & \multicolumn{3}{c}{5 Gdof} & & & & & &\\
\bottomrule
\end{tabular}
\end{center}
\end{table}

TODO

% \begin{table}[ht]%WAS h
% \caption{Single-node hardware specifications}
% \begin{tabular}{ccccccccc}
% \toprule
% Architecture & Model & Released & SIMD & Sockets & Cores & Base Freq. & Cache & DRAM\\
% & & & & & & (GHz) & (MB) & (GB)\\
% \midrule
% Barcelona & 8350B3 & 2008Q2 & SSE4a & 4 & 4 & 2.00 & 2+2+0.5 & 62\\
% Westmere & X5675 & 2011Q1 & SSE4.2 & 2 & 6 & 3.06 & 12 & 24\\
% Sandy Bridge & E5-2670 & 2012Q1 & AVX & 2 & 8 & 2.60 & 20+2+0.25 & 32-256\\
% Ivy Bridge & E5-2650V2 & 2013Q3 & AVX & 2 & 8 & 2.60 & 20 & 32-256\\
% Haswell & E5-2640V3 & 2014Q4 & AVX2 & 2 & 8 & 2.60 & 20 & 32-256\\
% Skylake A(B) & 6148 & 2017Q3 & AVX-512 & 2 & 20 & 2.40 & 27.5+20+1.25 & 96(375)\\
% Cascade Lake & 5118 & 2019Q2 & AVX-512 & 2 & 20 & 2.50 & 27.5+20+1.25 & 192\\
% \bottomrule
% \end{tabular}
% \end{table}

\begin{table}[ht]%WAS h
\caption{Single-node hardware specifications}
\begin{tabular}{ccccccccc}
\toprule
Architecture & SIMD & Cores & Base Freq. & Cache & DRAM\\
& & (total)  & (GHz) & (MB) & (GB)\\
\midrule
Barcelona &  SSE4a & 16 & 2.00 & 2+2+0.5 & 62\\
Westmere &  SSE4.2 & 12 & 3.06 & 12 & 24\\
Sandy Bridge & AVX & 16 & 2.60 & 20+2+0.25 & 32-256\\
Ivy Bridge & AVX & 16 & 2.60 & 20 & 32-256\\
Haswell  & AVX2 & 16 & 2.60 & 20 & 32-256\\
Skylake A(B) &  AVX-512 & 40 & 2.40 & 27.5+20+1.25 & 96(375)\\
Cascade Lake &  AVX-512 & 40 & 2.50 & 27.5+20+1.25 & 192\\
\bottomrule
\end{tabular}
\end{table}

\begin{table}[ht]%WAS h
\caption{Single-node hardware specifications}
\begin{tabular}{ccccccccc}
\toprule
Architecture & Model & Released & SIMD & Sockets & Cores \\
& & & & & \\
\midrule
Barcelona & 8350B3 & 2008Q2 & SSE4a & 4 & 4\\
Westmere & X5675 & 2011Q1 & SSE4.2 & 2 & 6 \\
Sandy Bridge & E5-2670 & 2012Q1 & AVX & 2 & 8 \\
Ivy Bridge & E5-2650V2 & 2013Q3 & AVX & 2 & 8 \\
Haswell & E5-2640V3 & 2014Q4 & AVX2 & 2 & 8 \\
Skylake A(B) & 6148 & 2017Q3 & AVX-512 & 2 & 20\\
Cascade Lake & 5118 & 2019Q2 & AVX-512 & 2 & 20\\
\bottomrule
\end{tabular}
\end{table}
TODO

\begin{figure}[ht]%WAS h
\begin{center}
\begin{minipage}[c]{0.48\linewidth}
\begin{center}
\includegraphics[width=\linewidth]{fig/femera-profile-nodes-tet10-elas-iso-gcc.eps}
%\caption{
Single-node performance profiles
%}
\end{center}
\end{minipage}
\begin{minipage}[c]{0.48\linewidth}
\begin{center}
\includegraphics[width=\linewidth]{fig/femera-profile-tet10-elas-iso-gcc.eps}
%\caption{
Per-core performance profiles
%}
\end{center}
\end{minipage}
\caption{Linear-elastic isotropic quadratic tetrahedral model performance profiles}
\label{fig:profile-tet10-core}
\end{center}
\end{figure}

Intel Advisor~\cite{ref:intel-advisor-2017} is a tool to collect detailed performance statistics while running an application. It also provides reference hardware bandwidth and computational specifications for the roofline
analyses discussed in section~\ref{subsec:efficient-small-tensor-eval}. 

\begin{figure}[ht]%WAS h
\begin{center}
\includegraphics[width=\linewidth]{fig/femera-core-time-50mdof-iso.eps}
\caption{Large quadratic tetrahedral model isotropic elastic processor kernel times using Intel CPUs}
\label{fig:time-nodes-bar}
\end{center}
\end{figure}

TODO

\subsection{Anisotropic and non-linear single-node performance}
\label{subsec:Linear-Elastic-Orthotropic}

TODO

\begin{figure}[ht]%WAS h
\begin{center}
\begin{minipage}[c]{1.0\linewidth}
\begin{center}
\includegraphics[width=\linewidth]{fig/femera-physics-time-6148.eps}
\end{center}
\end{minipage}
\begin{minipage}[c]{1.0\linewidth}
\begin{center}
\includegraphics[width=\linewidth]{fig/femera-physics-memory-6148.eps}
\end{center}
\end{minipage}
\caption{Large model kernel time and memory required for several types of physics evaluations (Skylake~A):
iteration time~(top) and memory required~(bottom).}
\label{fig:time-size-bar}
\end{center}
\end{figure}

TODO

\subsection{Computationally efficient small tensor operations}
\label{subsec:efficient-small-tensor-eval}

Both Intel version 2019.4.243 and GNU version 8.3.0 compiler suites can efficiently auto-vectorize all three PCG and the ISO F computational kernels in Algorithm~\ref{alg:iso-elastic-system-response} using double-precision~(DP) 128-bit streaming SIMD extensions~(SSE). Only the GNU compiler auto-vectorizes the ISO GH kernel efficiently with SSE instructions. 

Neither compiler can produce efficient 256-bit advanced vector extension~(AVX) or 512-bit~(AVX-512) instructions for the tensor physics kernels. However, significant performance improvements can be realized by specifying intrinsic AVX operations to implement the physics kernels. Memory-aligning and padding some of the data structures~(increasing their width from $n\times3$ to $n\times4$) improves performance. Although this strategy uses only three of the four available AVX DP SIMD lanes, and may effectively occupy less than half of dedicated AVX-512 hardware, section~\ref{sec:Examples} documents how well this strategy performs.

\begin{table}[ht]%WAS h
\caption{Typical preconditioned conjugate gradient iteration operations}
\label{tab:kern-pcg}
\begin{center}
\begin{tabular}{rrrr}
\toprule
Kernel & Computation & L1 Access & Arithmetic Intensity\\
&~(flop/dof) &~(byte/dof) &~(flop/byte)\\
\midrule
PCG A & 2 & 16 & 0.125\\
PCG B & 7 & 87 & 0.080\\
PCG C & 5 & 70 & 0.071\\
\cmidrule(lr){1-4}
Overall PCG & 14 & 173 & 0.081\\
\cmidrule(lr){1-4}
Other & – & 53 & –\\
\bottomrule
\end{tabular}
\end{center}
\end{table}

\emph{[Address why Kernels B and C FLOP counts are greater than what Algorithm~\ref{alg:halo-pcg} suggests.]}

TODO

\begin{table}[ht]%WAS h
\caption{Typical quadratic tetrahedron isotropic elastic response operations}
\label{tab:kern-iso-tet10}
\begin{tabular}{rrrrrrrr}
\toprule
& & & \multicolumn{2}{c}{Computation} & L1 Access & \multicolumn{2}{c}{Arithmetic Intensity}\\
& & &~(flop/dof) & Effective &~(byte/dof) & Measured & Effective\\
\cmidrule(lr){4-5}\cmidrule(lr){8-8}
Kernel Version & Vectorization & & Measured & & Measured &  & \\
\midrule
ISO GH 1 & Automatic* & SSE & 333 & & 1938 & 0.172 & \\
2 & Intrinsics & AVX & 440 & 333 & 799 & 0.550 & 0.417\\
\midrule
ISO S 1 & None & Scalar & 78 & & 1017 & 0.076 & \\
2 & Intrinsics** & AVX & 25 &  & 464 & 0.054 & \\
3 & Intrinsics & AVX2 & 58 &  & 8 & 7.250 & \\
\midrule
ISO F 1 & Automatic & SSE & 145 & & 581 & 0.250 & \\
2 & Intrinsics & AVX & 240 & 145 & 879 & 0.273 & 0.165\\
\midrule
Non-kernel & GH1 S1 F1 & SSE & - & & 446 & – & \\
& GH2 S2 F2 & AVX & 2 & & 837 & – & \\
& GH2 L3 F2 & AVX2 & 2 & & 845 & – & \\
\cmidrule(lr){1-8}
Overall Physics & Automatic & SSE & 556 &  & 3982 & 0.140 & \\
& Intrinsics & AVX & 707 & 556 & 2980 & 0.237 & 0.187\\
& Intrinsics & AVX2 & 740 & 556 & 2532 & 0.291 & 0.220\\
\bottomrule
\multicolumn{8}{l}{* The GNU compiler efficiently auto-vectorizes the ISO GH kernel for SSE.}\\
\multicolumn{8}{l}{** Partially vectorized}
\end{tabular}
\end{table}

TODO

\begin{table}[ht]%WAS h
\caption{Potential double-precision DRAM-limited isotropic elastic speedup}
\begin{tabular}{cccccccccc}
\toprule
Physics & & & & Operations & DRAM BW & Intensity & \multicolumn{2}{c}{Tensor Train Speedup}\\
Evaluation & & &~(flop/elem) &~(byte/elem) &~(flop/dof) &~(byte/dof) &~(flop/byte) & Physics & Iteration\\
\midrule
SpMV & P1 & & & 81/dof & 714 & 0.11 & 0.95 & 0.95\\
& P2 & & & 170/dof & 1447 & 0.12 & 7.1 & 4.9\\
& P3 & & & 285/dof & 2390 & 0.12 & 20.6 & 11.0\\
\midrule
LMS & P1 & & & & 1788 & 0.32 & 2.4 & 2.2\\
& P2 & & & & 1072 & 0.41 & 5.3 & 3.8\\
& P3 & & & & 1158 & 0.45 & 10.0 & 5.6\\
\midrule
Tensor Train & P1 & & & & 753 & 0.91 & 1 & 1\\
& P2 & & & & 203 & 3.21 & 1 & 1\\
& P3 & & & & 116 & 7.59 & 1 & 1\\
\bottomrule
\end{tabular}
\end{table}

TODO

\begin{figure}[ht]%WAS h
\begin{center}
\begin{minipage}[c]{0.48\linewidth}
\begin{center}
\includegraphics[width=\linewidth]{fig/femera-roofline-X5675-50mdof-tet10-iso.eps}
\caption{Computational kernels~(P2)}
\end{center}
\end{minipage}
\begin{minipage}[c]{0.48\linewidth}
\begin{center}
\includegraphics[width=\linewidth]{fig/femera-roofline-X5675-50mdof-pall-iso.eps}
\caption{Overall (P1, P2, P3)}
\end{center}
\end{minipage}
\caption{Typical auto-vectorized~(SSE) isotropic elastic large model performance roofline }
\label{fig:fig:roof-large-tet10-sse}
\end{center}
\end{figure}

TODO

\begin{figure}[ht]%WAS h
\begin{center}
\begin{minipage}[c]{0.48\linewidth}
\begin{center}
\includegraphics[width=\linewidth]{fig/femera-roofline-6148-50mdof-tet10-iso.eps}
\caption{Computational kernels~(P2)}
\end{center}
\end{minipage}
\begin{minipage}[c]{0.48\linewidth}
\begin{center}
\includegraphics[width=\linewidth]{fig/femera-roofline-6148-50mdof-pall-iso.eps}
\caption{Overall (P1, P2, P3)}
\end{center}
\end{minipage}
\caption{Typical intrinsics (AVX, AVX2) large model performance roofline}
\label{fig:roof-large-tet10-avx2}
\end{center}
\end{figure}

TODO

\begin{figure}[ht]%WAS h
\begin{center}
\begin{minipage}[c]{0.48\linewidth}
\begin{center}
\includegraphics[width=\linewidth]{fig/femera-roofline-6148-50mdof-tet10-lms.eps}
\caption{Computational kernels~(P2)}
\end{center}
\end{minipage}
\begin{minipage}[c]{0.48\linewidth}
\begin{center}
\includegraphics[width=\linewidth]{fig/femera-roofline-6148-50mdof-pall-lms.eps}
\caption{Overall (P1, P2, P3)}
\end{center}
\end{minipage}
\caption{Typical local matrix storage large model performance roofline}
\label{fig:roof-large-tet10-lms}
\end{center}
\end{figure}

TODO

\subsection{Domain over-decomposition cache blocking}
\label{subsec:Domain-Decomposition}

TODO

\begin{figure}[ht]%WAS h
\begin{center}
\begin{minipage}[c]{0.48\linewidth}
\begin{center}
\includegraphics[scale=0.5]{fig/femera-decomp-p2-iso-6148b.eps}
\caption{Partition performance}
\end{center}
\end{minipage}
\begin{minipage}[c]{0.48\linewidth}
\begin{center}
\includegraphics[scale=0.5]{fig/femera-part-whisker-p2-iso-6148b.eps}
\caption{Partition performance}
\end{center}
\end{minipage}
\caption{Partition Performance}
\label{fig:part-perf}
\end{center}
\end{figure}

TODO

\begin{table}[ht]%WAS h
\caption{Partitioning Parameters}
\begin{tabular}{ccccccc}
\toprule
Architecture & Element & Small & Large & & & \\
& Polynomial & (kdof) & (kdof) & (part) & (elem/part) & (kdof/part)\\
\midrule
Westmere & P1 & < & >719 & 108 & 13824 & 6.659\\
& P2 & < & >702 & 60 & 2926 & 11.707\\
& P3 & < & >1714 & 84 & 4503 & 20.415\\
\toprule
Sandy Bridge & P1 & < & >638 & 64 & 20704 & 9.973\\
& P2 & < & >3246 & 176 & 4610 & 18.444\\
& P3 & < & >2252 & 288 & 5751 & 78.211\\
\toprule
Ivy Bridge & P1 & < & >1205 & 112 & 22329 & 10.755\\
& P2 & < & >2458 & 96 & 4478 & 25.609\\
& P3 & < & >4673 & 272 & 1883 & 17.914\\
\toprule
Haswell & P1 & < & >681 & 64 & 22103 & 10.646\\
& P2 & < & >631 & 48 & 3287 & 13.152\\
& P3 & < & >1890 & 160 & 868 & 11.811\\
\toprule
Skylake A & P1 & < & >5603 & 240 & 48464 & 23.344\\
& P2 & < & >2696 & 80 & 8423 & 33.698\\
& P3 & < & >16837 & 320 & 3868 & 52.615\\
\toprule
Skylake B & P1 & < & >5216 & 240 & 45121 & 21.734\\
& P2 & < & >10318 & 240 & 10747 & 42.994\\
& P3 & < & >15434 & 240 & 4728 & 64.307\\
\toprule
Cascade Lake & P1 & < & > & & & \\
& P2 & < & > & & & \\
& P3 & < & > & & & \\
\bottomrule
\end{tabular}
\end{table}

TODO

\subsection{Nested parallel concurrent model solves}

TODO

\subsection{Strong and weak scaling}
\label{subsec:Strong-and-Weak}

TODO

\begin{figure}[ht]%WAS h
\begin{center}
\begin{minipage}[c]{0.48\linewidth}
\begin{center}
\includegraphics[width=\linewidth]{fig/femera-scale-500kdof-6148-iso.eps}
\caption{Strong and weak scaling to 1 Mdof}
\end{center}
\end{minipage}
\begin{minipage}[c]{0.48\linewidth}
\begin{center}
\includegraphics[width=\linewidth]{fig/femera-scale-50mdof-6148-iso.eps}
\caption{Strong and weak scaling to 100 Mdof}
\end{center}
\end{minipage}
\caption{Typical scaling of tensor train performance }
\label{fig:strong-weak}
\end{center}
\end{figure}

TODO

\section{Conclusions}
\label{sec:Conclusions}
%HJ 9. Reward ----------------------------------------------------------------
%      the hero experiences the consequences of surviving death

TODO

\section{Future Work}
\label{sec:Future-Work}

TODO

\subsection{multi-node}
\label{subsec:multi-node-mpi}
TODO

\subsection{Element and node sorting}
\label{subsec:Renumbering}

TODO

\subsection{Other preconditioners (multigrid, etc.)}
\label{subsec:Other-Preconditioners}

TODO

\subsection{Other elements (bricks, etc.)}
\label{subsec:Other-Elements-1}

TODO

\subsection{Mo physics! (phase field, etc.)}
\label{subsec:Mo-physics!}

TODO

\appendix


\section{Shape functions and integration rules}
\label{app:shapes-integration}

\section{Material property transformations}
\label{app:mtrl-xform}
%TODO compare Voigt matrix vs. tensor, construction and rotation

\section*{Notes}

TODO

% Draft be Claude 4.5 Sonnet 2026-01-21
\subsubsection{Billion-DOF Isotropic Validation}

To verify correctness at extreme scale, a nearly billion degree-of-freedom
(954M DOF) model was analyzed using 40,000 partitions on a single 40-core
node. The model consisted of a unit cube under uniform tension, for which
the analytical solution is uniform strain.

\textbf{Problem specification:}
\begin{itemize}
\item Geometry: Unit cube ${[0,1]}^3$
\item Material: Isotropic elastic ($E = 100 GPa, \nu = 0.3$)
\item Boundary conditions:
  \begin{itemize}
  \item x = 0: $u_x = 0$ (fixed)
  \item x = 1: $u_x = 0.001$ (prescribed displacement)
  \item Other faces: traction-free
  \end{itemize}
\item Mesh: 240M tetrahedral elements (P2), 318M nodes
\item Partitioning: 40,000 partitions (extreme over-decomposition)
\end{itemize}

\textbf{Analytical solution:}
\begin{itemize}
\item Uniform strain: $\epsilon_x = 0.001$
\item Uniform stress: $\sigma_x = E \cdot \epsilon_x = 100 MPa$
\item Lateral strain: $\epsilon_y = \epsilon_z = -\nu \cdot \epsilon_x = -0.0003$
\end{itemize}

\textbf{Results:}

Table~\ref{tab:billion-dof-validation} shows the solution accuracy after
convergence to a relative residual norm of $1.02 \times 10^{-82}$.

\begin{table}[ht]
\caption{Billion-DOF validation results}
\label{tab:billion-dof-validation}
\begin{tabular}{lccl}
\toprule
Quantity & Expected & Computed & Relative Error \\
\midrule
Applied stress (MPa) & 100.0 & 100.0 & - \\
Extension strain & 0.001000 & 0.001000 & $1.33 \times 10^{-14}$ \\
Back-calculated E (GPa) & 100.0 & 100.0 & $1.33 \times 10^{-14}$ \\
\midrule
\multicolumn{4}{l}{\textbf{Displacement field errors:}} \\
Maximum (any component) & 0 & $5.63 \times 10^{-13}$ & - \\
Average magnitude & 0 & $1.05 \times 10^{-13}$ & - \\
\bottomrule
\end{tabular}
\end{table}

The solution reproduces the analytical uniform strain state to within
machine precision (errors < $10^{-12}$), demonstrating:
\begin{itemize}
\item Correct element formulation and integration
\item Accurate material property application  
\item Proper stress-strain relationships
\item No accumulation of numerical errors over 954M DOF system
\item Correct parallel assembly with 40,000 partitions
\end{itemize}

The problem required 100,000 PCG iterations to reach machine precision,
solving in 52.2 hours wall-clock time (2.125 days) on a single 40-core
Xeon 6148 (Skylake) node. The aggregate physics kernel performance was
454.65 GFLOP/s with an arithmetic intensity of 1.24 FLOP/byte, consistent
with DRAM bandwidth-limited performance predicted by roofline analysis
%(Section~\ref{sec:roofline}).

Parallel efficiency was excellent despite extreme over-decomposition
(40,000 partitions on 40 threads, 1000:1 partition-to-thread ratio):
\begin{itemize}
\item Load balance: 101.69\% (1.69\% imbalance)
\item Synchronization overhead: 4.90\% of total time
\item Useful work: 95.10\% of total time
\end{itemize}

This test validates both correctness and scalability of the tensor train
element-by-element approach for problems approaching one billion degrees
of freedom on single-node architectures.

\bibliographystyle{plain}
\bibliography{femera}{}

\end{document}

%%% Local Variables: 
%%% mode: latex
%%% TeX-master: t
%%% End: 
